\documentclass{beamer}
%\documentclass[handout]{beamer}
\mode<presentation>
{
%  \usetheme{Darmstadt}%default Warsaw
  \setbeamercovered{transparent}
}

\setbeamertemplate{navigation symbols}{} 

%\useinnertheme[shadow=true]{rounded}

\setbeamercolor{block title}{use=structure,fg=structure.fg,bg=structure.fg!20!bg}
\setbeamercolor{block body}{parent=normal text,use=block title,bg=block title.bg!50!bg}

\usepackage{beamerthemesplit}
\usepackage{epsfig,graphics}
\usepackage{amsmath,amsfonts,amssymb,color,array,subfigure,rotating,texdraw} 
\usepackage{multimedia}
\usepackage{psfrag}
%\usepackage[dvips]{color}
\usepackage[polish]{babel}

\newcommand{\bDiamond}{\mathbin{\Diamond}}
\newcommand{\bLozenge}{\mathbin{\blacklozenge}}

\makeatletter
\newcommand\bigDiamond{\mathop{\mathpalette\bigDi@mond\relax}}
\newcommand\bigDi@mond[2]{%
  \vcenter{\hbox{\m@th
    \scalebox{\ifx#1\displaystyle 2\else1.2\fi}{$#1\Diamond$}%
  }}%
}
\newcommand\bigLozenge{\mathop{\mathpalette\bigL@zenge\relax}}
\newcommand\bigL@zenge[2]{%
  \vcenter{\hbox{\m@th
    \scalebox{\ifx#1\displaystyle 2\else1.2\fi}{$#1\blacklozenge$}%
  }}%
}
\makeatother


\definecolor{greenn}{rgb}{0.1,0.4,0.1}

\def\Dpartial#1#2{ {\partial #1 \over \partial #2} }
\def\Dpartialmix#1#2#3{ {\partial^2 #1 \over \partial #2\,\partial #3} }
\def\Dpartialn#1#2#3{ {\partial^{#3} #1 \over \partial #2^{#3}} }
\def\Txtint{ {\textstyle\int} }
\def\Txtiint{ {\textstyle\iint} }
\def\Bmp#1{ \begin{minipage}{#1} }
\def\Emp{ \end{minipage} }
\def\Bmpc#1{ \begin{minipage}[c]{#1} }
\def\Bmpt#1{ \begin{minipage}[t]{#1} }
\def\Bmpb#1{ \begin{minipage}[b]{#1} }
\newcommand{\stitle}[1]{\centering{{\sc \Large #1 }}\vskip .5cm}
\def\H{{\mathcal{H}}}
\def\J{{\mathcal{J}}}

\def\A{{\mathcal{A}}}
\def\B{{\mathcal{B}}}
\def\E{{\mathcal{E}}}
\def\K{{\mathcal{K}}}
\def\L{{\mathcal{L}}}
\def\M{{\mathcal{M}}}
\def\O{{\mathcal{O}}}
\def\F{{\mathcal{F}}}
\def\N{{\mathcal{N}}}
\def\T{{\mathcal{T}}}
\def\P{{\mathcal{P}}}
\def\R{{\mathcal{R}}}
\def\G{{\mathcal{G}}}
\def\K{{\mathcal{K}}}
\def\C{{\mathcal{C}}}
\def\U{{\mathcal{U}}}
\def\S{{\mathcal{S}}}
\def\DD{{\mathcal{D}}}
\def\QQ{{\mathcal{Q}}}
\def\NS{{\mathcal{NS}}}
\def\l{{\ell}}

\def\BB{{\mathbb{B}}}
\def\KK{{\mathbb{K}}}
\def\RR{{\mathbb{R}}}
\def\TT{{\mathbb{T}}}
\def\ZZ{{\mathbb{Z}}}
\def\CC{{\mathbb{C}}}

\def\tu{{\tilde{u}}}
\def\ta{{\tilde{a}}}
\def\tb{{\tilde{b}}}

\def\magenta#1{{\color{magenta}{ #1 }}}
\def\red#1{{\color{red}{ #1 }}}
\def\blue#1{{\color{blue}{ #1 }}}
\def\green#1{{\color{green}{ #1 }}}
\def\white#1{{\color{white}{ #1 }}}
\def\black#1{{\color{black}{ #1 }}}

\def\bC{{\bf C}}
\def\bG{{\bf G}}
\def\bZ{{\bf Z}}

\newcommand{\bds}{\boldsymbol}
\newcommand{\mbb}[1]{\mathbb{#1}}
\newcommand{\mc}[1]{\mathcal{#1}}

\def\P{{\mathcal{P}}}
\def\b{{\bf b}}
\def\f{{\bf f}}
\def\q{{\bf q}}
\def\r{{\bf r}}
\def\v{{\bf v}}
\def\z{{\bf z}}
\def\y{{\bf y}}
\def\k{{\bf k}}
\def\x{{\bf x}}
\def\m{{\bf m}}
\def\l{{\bf l}}
\def\n{{\bf n}}
\def\e{{\bf e}}
\def\j{{\bf j}}
\def\D{{\bf D}}
\def\X{{\bf X}}
\def\Q{{\bf Q}}
\def\g{{\bf g}}
\def\u{{\bf u}}
\def\f{{\bf f}}
\def\y{{\bf y}}
\def\0{{\mathbf{0}}}
\def\V{{\mathbf{V}}}
\def\W{{\mathbf{W}}}
\def\Z{{\mathbf{Z}}}

\def\tT{{\tilde{\mathcal{T}}}}
\def\tZ{{\tilde{\mathcal{Z}}}}

\newcommand{\tuE}{\widetilde{\mathbf{u}}_{\E_0}}
\newcommand{\tuET}{\widetilde{\u}_{0;\E_0,T}}
\newcommand{\uET}{\u_{0;\E_0,T}}
\newcommand{\ET}{\E_T(\u_0)}
\def\tTE{\widetilde{T}_{\E_0}}

\def\btau{{\bf \tau}}
\def\hu{{\hat{u}}}
\def\bnabla{\boldsymbol{\nabla}}
\def\bphi{\boldsymbol{\phi}}
\def\bDelta{\boldsymbol{\Delta}}
\def\bPhi{\boldsymbol{\Phi}}
\def\bxi{\boldsymbol{\xi}}
\def\etab{\boldsymbol{\eta}}
\def\control{{\phi}}
\def\controlpert{{\phi'}}
\def\bsigma{\boldsymbol{\sigma}}
\def\bomega{\boldsymbol{\omega}}
\def\teta{\widetilde{\boldsymbol{\eta}}}

\newcommand{\argmin}{\operatorname{argmin}}
\newcommand{\argmax}{\operatorname{argmax}}
\newcommand{\sgn}{\operatorname{sgn}}
\newcommand{\sign}{\operatorname{sign}}
\newcommand{\rank}{\operatorname{rank}}
\newcommand{\Ext}{\operatorname{Ext}}
\newcommand{\id}{\,\mathrm{d}}
\newcommand{\Id}{\operatorname{Id}}
\newcommand{\Ker}{\operatorname{Ker}}
\newcommand{\kvec}{\mathbf{k}}

\newcommand{\sigdis}{\sigma_{\text{disc}}}
\newcommand{\sigess}{\sigma_{\text{ess}}}
\def\sigp{\sigma_{\text{p}}}
\def\mumax{\mu_{\text{max}}}

\long\def\symbolfootnote[#1]#2{\begingroup%
\def\thefootnote{\fnsymbol{footnote}}\footnote[#1]{#2}\endgroup} 

%%%%%%%%%%%%%%%%%%%%%%%%%%%%%%%%%%%%%%%%%%%%%%%%%%%%%%%%%%%%%%%%%%%%%%%%%%
%%%%%%%%%%%%%%%%%%%%%%%%%%%%%%%%%%%%%%%%%%%%%%%%%%%%%%%%%%%%%%%%%%%%%%%%%%
%%%%%%%%%%%%%%%%%%%%%%%%%%%%%%%%%%%%%%%%%%%%%%%%%%%%%%%%%%%%%%%%%%%%%%%%%%

\title[Search for Singularities and Instabilities in Euler Flows]
{Systematic Search for Singularities and Instabilities in Euler Flows}
%\subtitle{Extreme Vortex States and \\ the Hydrodynamic Blow-Up Problem}
\author[B.~Protas]{Bartosz Protas}
\institute{
Department of Mathematics \& Statistics \\
McMaster University, Hamilton, Ontario, Canada \\
URL:   \href{https://ms.mcmaster.ca/~bprotas/}{\tt http://www.math.mcmaster.ca/bprotas} \\ \vspace*{0.3cm}

%\footnotesize
Funded by NSERC (Canada), Computing Time Provided by {\sc Compute Canada}}
%\vspace*{-0.2cm}

\date{\scriptsize  GFD Program @ Woods Hole Oceanographic Institution   \\ \medskip
July 10, 2025  \\ \smallskip}


\begin{document}

\frame{\titlepage}


%%%%%%%%%%%%%%%%%%%%%%%%%%%%%%%%%%%%%%%%%%%%%%%%%%%%%%%%%%%%%%%%%%%%%%%%%%%
%%%%%%%%%%%%%%%%%%%%%%%%%%%%%%%%%%%%%%%%%%%%%%%%%%%%%%%%%%%%%%%%%%%%%%%%%%%
%%%%%%%%%%%%%%%%%%%%%%%%%%%%%%%%%%%%%%%%%%%%%%%%%%%%%%%%%%%%%%%%%%%%%%%%%%%

 %%\section[Agenda]{}
 %\frame{
 %\frametitle{In Memoriam: \ Charles R.~Doering (1956--2021)}

 %\centering
 %\includegraphics[width=0.5\textwidth]{Figs9/doering.jpg}  

 %}


%%%%%%%%%%%%%%%%%%%%%%%%%%%%%%%%%%%%%%%%%%%%%%%%%%%%%%%%%%%%%%%%%%%%%%%%%%
%%%%%%%%%%%%%%%%%%%%%%%%%%%%%%%%%%%%%%%%%%%%%%%%%%%%%%%%%%%%%%%%%%%%%%%%%%
%%%%%%%%%%%%%%%%%%%%%%%%%%%%%%%%%%%%%%%%%%%%%%%%%%%%%%%%%%%%%%%%%%%%%%%%%%

%\section[Agenda]{}
\frame{
\frametitle{Agenda}
\tableofcontents
}

%%%%%%%%%%%%%%%%%%%%%%%%%%%%%%%%%%%%%%%%%%%%%%%%%%%%%%%%%%%%%%%%%%%%%%%%%%
%%%%%%%%%%%%%%%%%%%%%%%%%%%%%%%%%%%%%%%%%%%%%%%%%%%%%%%%%%%%%%%%%%%%%%%%%%
%%%%%%%%%%%%%%%%%%%%%%%%%%%%%%%%%%%%%%%%%%%%%%%%%%%%%%%%%%%%%%%%%%%%%%%%%%
\section{Systematic Search for Singularities via PDE Optimization}
\frame{
\centering

\LARGE
{\sc PART I: \\ \medskip Systematic Search for Singularities via PDE Optimization}

\normalsize
\bigskip\bigskip
Joint work with:
\begin{itemize}
\item Di Kang (former post-doc, now in industry)
\medskip
\item Xinyu Zhao (post-doc, Assistant Professor at the New Jersey Institute of Technology as of September 2024)
\end{itemize}
}

%%%%%%%%%%%%%%%%%%%%%%%%%%%%%%%%%%%%%%%%%%%%%%%%%%%%%%%%%%%%%%%%%%%%%%%%%%
%%%%%%%%%%%%%%%%%%%%%%%%%%%%%%%%%%%%%%%%%%%%%%%%%%%%%%%%%%%%%%%%%%%%%%%%%%
%%%%%%%%%%%%%%%%%%%%%%%%%%%%%%%%%%%%%%%%%%%%%%%%%%%%%%%%%%%%%%%%%%%%%%%%%%

\frame{
%\frametitle{References}
\centering
\medskip
\hspace*{-2.5cm} \includegraphics[width=0.575\textwidth]{Figs9/kyp19.png}
\medskip
\hspace*{2.5cm} \includegraphics[width=0.6\textwidth]{Figs9/kp22.png}
\medskip
\hspace*{-2.5cm} \includegraphics[width=0.6\textwidth]{Figs9/zp23a_v2.jpg}

\bigskip\smallskip
}


%%%%%%%%%%%%%%%%%%%%%%%%%%%%%%%%%%%%%%%%%%%%%%%%%%%%%%%%%%%%%%%%%%%%%%%%%%
%%%%%%%%%%%%%%%%%%%%%%%%%%%%%%%%%%%%%%%%%%%%%%%%%%%%%%%%%%%%%%%%%%%%%%%%%%
%%%%%%%%%%%%%%%%%%%%%%%%%%%%%%%%%%%%%%%%%%%%%%%%%%%%%%%%%%%%%%%%%%%%%%%%%%

%\section{Sharpness of Estimates as Optimization Problem}
\frame
{
\hspace*{-1.0cm}
\Bmp{12cm}
\raggedright

\begin{itemize}
\item<1->
Navier-Stokes system ($\Omega = [0,L]^d$, $d=2,3$)
%\footnotesize
\small
\begin{equation*}
\left\{
\begin{aligned}
&\Dpartial{\u}{t}+({\u}\cdot{\bnabla}){\u}
+{\bnabla}{p}-\nu\Delta{{\u}}=\0,  && \text{in} \ \Omega\times(0,T] \\ 
&{\bnabla}\cdot{\u}=0,                     && \text{in} \ \Omega\times(0,T] \\ 
&\u = \u_0                 &&  \text{in} \ \Omega \ \text{at} \ t=0 \\
&\textrm{Periodic Boundary Conditions}                  &&  \text{on} \ \Gamma\times(0,T] 
\end{aligned}
%\label{eq:NS}
\right.
\end{equation*}
\medskip

\item<2->
{\bf The Big Question:} \\ \smallskip
{\blue{\em   Given a smooth initial condition $\u_0$, does the Navier-Stokes
  system always admit smooth solutions $\u(t)$ for arbitrarily long
  times $t$? }} \\ \smallskip
\visible<3->{\blue{(solutions which are not ``smooth'' are not physically meaningful ...)}}


\bigskip

\item<4->
One of the Clay Institute ``Millennium Problems'' (\$ 1M prize!) \\
\smallskip
\scriptsize
\qquad
{\tt{http://www.claymath.org/millennium/Navier-Stokes\_Equations}}

\end{itemize}
\Emp
}

%%%%%%%%%%%%%%%%%%%%%%%%%%%%%%%%%%%%%%%%%%%%%%%%%%%%%%%%%%%%%%%%%%%%%%%%%%
%%%%%%%%%%%%%%%%%%%%%%%%%%%%%%%%%%%%%%%%%%%%%%%%%%%%%%%%%%%%%%%%%%%%%%%%%%
%%%%%%%%%%%%%%%%%%%%%%%%%%%%%%%%%%%%%%%%%%%%%%%%%%%%%%%%%%%%%%%%%%%%%%%%%%

\frame
{
\hspace*{-1.0cm}
\Bmp{12cm}
\raggedright

\begin{itemize}
\item<1->
What could go wrong with solutions to the Navier-Stokes equation?
\medskip

\item<2->
Consider its vorticity formulation ($\bomega = \bnabla \times \u$)
%\footnotesize
\small
\begin{equation*}
\Dpartial{\bomega}{t}+({\u}\cdot{\bnabla}){\bomega} = \frac{d \bomega}{dt} 
= \nu\Delta{{\bomega}} + \hspace*{-0.3cm} \underbrace{(\bnabla\u)\,\bomega}_{\text{``vortex stretching''}} 
%= \nu\Delta{{\omega}} + (\bnabla\u)\,\bomega 
\end{equation*}
\smallskip

\normalsize
\item<3->
Velocity $\u$ is obtained from vorticity using the Biot-Savart kernel $\bG$
\small
\begin{equation*}
\u = \bnabla \Delta^{-1} \bomega = \int_{\Omega} \bG(\cdot,\x')\bomega(\x')\, d\x' = \bG * \bomega
\end{equation*}

\normalsize
\item<4->
The vorticity equation has a quadratic source term (assume $\nu = 0$)
\small
\begin{equation*}
\frac{d \bomega}{dt} 
= \blue{\left[ \bnabla (\bG * \bomega) \right] \,\bomega}
\end{equation*}

\normalsize
\item<5->
What could this imply?


\end{itemize}
\Emp
}


%%%%%%%%%%%%%%%%%%%%%%%%%%%%%%%%%%%%%%%%%%%%%%%%%%%%%%%%%%%%%%%%%%%%%%%%%%
%%%%%%%%%%%%%%%%%%%%%%%%%%%%%%%%%%%%%%%%%%%%%%%%%%%%%%%%%%%%%%%%%%%%%%%%%%
%%%%%%%%%%%%%%%%%%%%%%%%%%%%%%%%%%%%%%%%%%%%%%%%%%%%%%%%%%%%%%%%%%%%%%%%%%

\frame
{
\hspace*{-1.0cm}
\Bmp{12cm}
\raggedright

\begin{itemize}
\item<1->
Consider a (very) simple ODE model problem
\begin{equation*}
\frac{d y}{dt}  = y^2, \ y(0) = y_0 \quad \text{with solution} \quad
\blue{y(t) = \frac{y_0}{1 - y_0 \, t}}
\end{equation*}
\medskip 

\item<2->
The solution $y(t)$ becomes unbounded \\ 
(``blows up'') as $t \rightarrow \blue{t_0 = \frac{1}{y_0}}$ \\
\mbox{
\hspace*{6.75cm}
\Bmp{0.5\textwidth}
\vspace*{-1.25cm}
\visible<2->{\includegraphics[width=0.6\textwidth]{Figs9/y2_blowup_v1}}
\Emp
}
\medskip 

\item<3->
The equation is not satisfied at $t=t_0$ and \\ the solution is not defined for $t \ge t_0$

\end{itemize}

\Emp
}


%%%%%%%%%%%%%%%%%%%%%%%%%%%%%%%%%%%%%%%%%%%%%%%%%%%%%%%%%%%%%%%%%%%%%%%%%%
%%%%%%%%%%%%%%%%%%%%%%%%%%%%%%%%%%%%%%%%%%%%%%%%%%%%%%%%%%%%%%%%%%%%%%%%%%
%%%%%%%%%%%%%%%%%%%%%%%%%%%%%%%%%%%%%%%%%%%%%%%%%%%%%%%%%%%%%%%%%%%%%%%%%%

\frame
{
\hspace*{-1.0cm}
\Bmp{12cm}
\raggedright

\begin{itemize}
\item<1->
Another simple model problem --- inviscid Burgers equation
\small
\begin{equation*}
\begin{alignedat}{2}
\Dpartial{u}{t} + u \Dpartial{u}{x} & = 0& \quad &\text{for} \ t >0, \ x \in \RR, \\
u(0,x) & = \phi(x) & & \text{for} \ x \in \RR  
\end{alignedat}
\end{equation*}
with (implicit) solution \blue{$u(t,x) = \phi(x - u(t,x)\, t)$}
\medskip 

\item<2->
The solution $u(t,x)$ develops a shock and becomes non-differentiable 
at time  $t \rightarrow \blue{t_0 = \frac{-1}{\min_x \frac{d\phi(x)}{dx}}}$ \\
\mbox{
\hspace*{5.5cm}
\Bmp{0.5\textwidth}
\vspace*{-0.75cm}
\visible<2->{\includegraphics[width=0.7\textwidth]{Figs9/burgers_blowup_v2}}
\Emp
}
%\medskip 

\item<3->
The equation is not satisfied at $t=t_0$ and the solution is not defined (in the ``classical'' sense) for $t \ge t_0$
\begin{itemize}
\item<4-> it may be however defined for $t \ge t_0$ in a ``weak'' (integral) sense
\end{itemize}

\end{itemize}

\Emp
}

%%%%%%%%%%%%%%%%%%%%%%%%%%%%%%%%%%%%%%%%%%%%%%%%%%%%%%%%%%%%%%%%%%%%%%%%%%
%%%%%%%%%%%%%%%%%%%%%%%%%%%%%%%%%%%%%%%%%%%%%%%%%%%%%%%%%%%%%%%%%%%%%%%%%%
%%%%%%%%%%%%%%%%%%%%%%%%%%%%%%%%%%%%%%%%%%%%%%%%%%%%%%%%%%%%%%%%%%%%%%%%%%

\frame
{
\hspace*{-1.0cm}
\Bmp{12cm}
\raggedright

\begin{itemize}
\item<1->
Can such singular behavior arise in the Navier-Stokes system \\ in finite time?
\medskip 

\item<2->
Who cares?
\begin{itemize}
\item<3-> Well, if its solutions can become singular, then the
  Navier-Stokes system is not a correct model for viscous
  incompressible fluids and must be amended (by modifying the viscous
  terms)
\end{itemize}
\medskip 

\item<4->
2D Case
\begin{itemize}
\item<4->
Existence theory complete --- smooth and unique solutions exist for
arbitrary times and arbitrarily large data
\end{itemize}
\medskip 

\item<5->
3D Case
\begin{itemize}
\item<5->
Weak solutions (possibly nonsmooth) exist for arbitrary times
\smallskip 

\item<6->
Classical (smooth) solutions exist for {\it finite} times only 
\smallskip 

\item<7->
Possibility of ``blow-up'' (finite-time singularity formation)
\end{itemize}


\end{itemize}

\Emp
}




%%%%%%%%%%%%%%%%%%%%%%%%%%%%%%%%%%%%%%%%%%%%%%%%%%%%%%%%%%%%%%%%%%%%%%%%%%
%%%%%%%%%%%%%%%%%%%%%%%%%%%%%%%%%%%%%%%%%%%%%%%%%%%%%%%%%%%%%%%%%%%%%%%%%%
%%%%%%%%%%%%%%%%%%%%%%%%%%%%%%%%%%%%%%%%%%%%%%%%%%%%%%%%%%%%%%%%%%%%%%%%%%

\subsection{Conditional Regularity Results}
%\subsection{Bounds on Rate of Growth of Enstrophy}
\frame
{
\frametitle{The Enstrophy Condition}

\begin{itemize}
\item<1->
A Key Quantity --- Enstrophy
\begin{equation*}
\E(t) \triangleq \int_{\Omega} |\bnabla\times\u|^2\, d\Omega \qquad (= \| \bnabla\u \|_2^2 )
\end{equation*}

\item<2-> 
Smoothness of Solutions $\Longleftrightarrow$ Bounded Enstrophy \\ (Foias \& Temam, 1989)
\mbox{
\Bmp{0.4\textwidth}
\begin{equation*}
\max_{t \in [0,T]} \E(t) < \infty \quad ???
\end{equation*}
\Emp\qquad
\Bmp{0.5\textwidth}
\vspace*{-0.5cm}
\visible<2->{\includegraphics[width=0.6\textwidth]{Figs9/Et_blowup_01}}
\Emp
}

\item<3->
Can estimate $\frac{d\E(t)}{dt}$ using the momentum equation, Sobolev's
embeddings, Young and Cauchy-Schwartz inequalities, ...
\begin{itemize}
\item<4->
{\sc Remark:} incompressibility not used in these estimates ....
\end{itemize}


\end{itemize}
}


%%%%%%%%%%%%%%%%%%%%%%%%%%%%%%%%%%%%%%%%%%%%%%%%%%%%%%%%%%%%%%%%%%%%%%%%%%
%%%%%%%%%%%%%%%%%%%%%%%%%%%%%%%%%%%%%%%%%%%%%%%%%%%%%%%%%%%%%%%%%%%%%%%%%%
%%%%%%%%%%%%%%%%%%%%%%%%%%%%%%%%%%%%%%%%%%%%%%%%%%%%%%%%%%%%%%%%%%%%%%%%%%

\frame
{
%\frametitle{What is known? --- Available Estimates}

\begin{itemize}
\item<1->
Bounds on the rate of growth of enstrophy --- general form
\begin{equation*}
\frac{d\E}{dt} < C \, \E^{\alpha}, \quad C > 0, \quad \blue{\alpha = \alpha(d) > 0}
\end{equation*}
\medskip

\item<2->
Energy equation ($\K(t) \triangleq \int_{\Omega} \u^2\, d\Omega $)
\begin{align*}
\frac{d\K}{dt} &= - 2 \nu \E \\
\visible<3->{
\K(t) - \K(0) &= - 2 \nu \int_0^t \E(\tau)\, d\tau 
\quad \Longrightarrow \quad \blue{\int_0^t \E(\tau)\, d\tau \le \frac{1}{2\nu}\K_0}}
\end{align*}
\medskip

\item<4->
When $\alpha <= 2$, by Gr\"onwall's inequality: 
$\E(t) \le \E_0 \exp\left[\frac{C \K_0}{2\nu}\right]$ 
\begin{itemize}
\item<5->[] $\Longrightarrow$ Enstrophy bounded for {\em all} times
\end{itemize}
\medskip

\item<6->
When \red{$\alpha > 2$}, no finite a priori bound on enstrophy ...


\end{itemize}
}


%%%%%%%%%%%%%%%%%%%%%%%%%%%%%%%%%%%%%%%%%%%%%%%%%%%%%%%%%%%%%%%%%%%%%%%%%%
%%%%%%%%%%%%%%%%%%%%%%%%%%%%%%%%%%%%%%%%%%%%%%%%%%%%%%%%%%%%%%%%%%%%%%%%%%
%%%%%%%%%%%%%%%%%%%%%%%%%%%%%%%%%%%%%%%%%%%%%%%%%%%%%%%%%%%%%%%%%%%%%%%%%%



\frame
{
%\frametitle{What is known? --- Available Estimates}

\begin{itemize}
\item<1->
2D Case:
\begin{equation*}
\frac{d\E(t)}{dt} \le \frac{C^2}{\nu} \E(t)^2
\end{equation*}
\begin{itemize}
\item<2->
Gr\"onwall's lemma and energy equation yield $\forall_t$ $\E(t) < \infty$
\smallskip
\item<3->
smooth solutions exist for all times
\end{itemize}

\bigskip
\item<4->
3D Case:
\begin{equation*}
\frac{d\E(t)}{dt} \le \frac{27 C^2}{128 \nu^3} \E(t)^3
\end{equation*}
\begin{itemize}
\item<5->
upper bound on $\E(t)$ blows up in finite time
\begin{equation*}
\E(t) \le \frac{\E(0)}{\sqrt{1 - 4 \frac{C \E(0)^2}{\nu^3} t}}
\end{equation*}
\item<6->
singularity in finite time cannot be ruled out!
\end{itemize}

\end{itemize}
}

%%%%%%%%%%%%%%%%%%%%%%%%%%%%%%%%%%%%%%%%%%%%%%%%%%%%%%%%%%%%%%%%%%%%%%%%%%
%%%%%%%%%%%%%%%%%%%%%%%%%%%%%%%%%%%%%%%%%%%%%%%%%%%%%%%%%%%%%%%%%%%%%%%%%%
%%%%%%%%%%%%%%%%%%%%%%%%%%%%%%%%%%%%%%%%%%%%%%%%%%%%%%%%%%%%%%%%%%%%%%%%%%


%\subsection{The Ladyzhenskaya-Prodi-Serrin (LPS) Conditions}
\frame
{
%\frametitle{What is known? --- Available Estimates}
\hspace*{-1.0cm}
\Bmp{12cm}
\raggedright


\begin{itemize}
\item<1-> \blue{\sc The Ladyzhenskaya-Prodi-Serrin (LPS) conditions}
  --- \\ the solution $\u(t)$ is smooth and satisfies the
  Navier-Stokes system in the classical sense provided that
\begin{equation*}
\blue{\u \in L^p([0,T];L^q(\Omega)), \quad 2/p+3/q = 1, \quad q > 3}
%\label{eq:LPS}
\end{equation*}

\item<2->  Thus, should a singularity form 
at some finite time $0 < t_0 < \infty$, then necessarily \\ \vspace*{-0.3cm}
\begin{equation*}
\red{\lim_{t \rightarrow t_0} \int_0^t \| \u(\tau) \|_{L^q(\Omega)}^p \, d\tau \rightarrow \infty, \quad 2/p+3/q = 1, \quad q > 3,} \vspace*{-0.2cm}
%\label{eq:LPSblowup}
\end{equation*}
where \quad $\|\u(t))\|_{L^q(\Omega)}  := \left( \int_\Omega |\u(t,\x)|^q \,d\x \right)^{\frac{1}{q}}$ 

\bigskip
\item<3->
In the  limiting case with $q = 3$, the corresponding condition for regularity is
\begin{equation*}
\blue{\u \in L^{\infty}([0,T];L^3(\Omega))}
%\label{eq:LPS3}
\end{equation*}

%\medskip
%\item<4-> {\blue{\sc Question:}} How to look for initial data $\u_0$
%  which may lead to singularity in finite time?

\end{itemize}

\Emp
}


%%%%%%%%%%%%%%%%%%%%%%%%%%%%%%%%%%%%%%%%%%%%%%%%%%%%%%%%%%%%%%%%%%%%%%%%%%
%%%%%%%%%%%%%%%%%%%%%%%%%%%%%%%%%%%%%%%%%%%%%%%%%%%%%%%%%%%%%%%%%%%%%%%%%%
%%%%%%%%%%%%%%%%%%%%%%%%%%%%%%%%%%%%%%%%%%%%%%%%%%%%%%%%%%%%%%%%%%%%%%%%%%


%\subsection{The Ladyzhenskaya-Prodi-Serrin (LPS) Conditions}
\frame
{
\frametitle{On the Nature of Possible Blow-up}
\hspace*{-1.0cm}
\Bmp{12cm}
\raggedright


\begin{itemize}
\item<1->
As the hypothetical blow-up time $t_0$ is approached ...

\begin{itemize}
  \item<2->[$\bullet$]
  \begin{equation*}
    \lim_{t \rightarrow t_0} \E(t) = \infty  \qquad \hfill
    \text{however} \hfill 
    \qquad \red{\int_0^{t_0} \E(\tau) \, d\tau < \infty}
%\label{eq:LPSblowup}
  \end{equation*}
\bigskip

\item<3->[$\bullet$]
  \begin{equation*}
    \lim_{t \rightarrow t_0} \int_0^t \| \u(\tau) \|_{L^q(\Omega)}^p \, d\tau = \infty, \quad 2/p+3/q = 1, \quad q > 3,
%\label{eq:LPSblowup}
  \end{equation*}
  however
  \begin{equation*}
\red{\int_0^{t_0} \| \u(\tau) \|_{L^q(\Omega)}^{\frac{4q}{3(q-2)}} \,
  d\tau < \infty, \qquad 2 \le q \le 6}
\end{equation*}
\end{itemize}

\bigskip
\item<4-> Thus, the blow-up, should it occur, must be very gentle ...

\end{itemize}

\Emp
}


%%%%%%%%%%%%%%%%%%%%%%%%%%%%%%%%%%%%%%%%%%%%%%%%%%%%%%%%%%%%%%%%%%%%%%%%%%
%%%%%%%%%%%%%%%%%%%%%%%%%%%%%%%%%%%%%%%%%%%%%%%%%%%%%%%%%%%%%%%%%%%%%%%%%%
%%%%%%%%%%%%%%%%%%%%%%%%%%%%%%%%%%%%%%%%%%%%%%%%%%%%%%%%%%%%%%%%%%%%%%%%%%

%\subsection{Conditional Regularity Results for the Euler System}
\frame
{
%\frametitle{What is known? --- Available Estimates}
\hspace*{-0.8cm}
\Bmp{12cm}
\raggedright

\begin{itemize}
\item<1->
Euler system (obtained by setting $\nu = 0$ in the Navier-Stokes system)
%\footnotesize
\small
\begin{equation*}
\left\{
\begin{aligned}
&\Dpartial{\u}{t}+({\u}\cdot{\bnabla}){\u}
+{\bnabla}{p}=\0,  && \text{in} \ \Omega\times(0,T] \\ 
&{\bnabla}\cdot{\u}=0,                     && \text{in} \ \Omega\times(0,T] \\ 
&\u = \u_0                 &&  \text{in} \ \Omega \ \text{at} \ t=0 \\
&\textrm{Periodic Boundary Conditions}                  &&  \text{on} \ \Gamma\times(0,T] 
\end{aligned}
%\label{eq:NS}
\right.
\end{equation*}
%\medskip

\item<3->
Local well-posedness of classical solutions in Sobolev spaces \\ \vspace*{-0.2cm}
\centering \footnotesize
\begin{theorem}[Kato, 1972]
  If $\u_0 \in H^m (\TT^3)$ or $H^m (\RR^3)$ for some $m > 5/2$ and
  satisfies $\bnabla \cdot \u_0 = 0$, then there exists a time $T>0$
  such that the Euler system has a unique solution $\u = \u(t,\x)$ with the initial
  condition $\u_0$ and the solution satisfies $\u \in C([0,
  T]; H^m)\bigcap C^1([0,T]; H^{m-1})$.
\end{theorem}
\smallskip

\item<4->
Global well-posedness remains open --- does there exist a smooth initial condition $\u_0 \in H^m$ for $m>5/2$, such that \\ \vspace*{-0.3cm}
\begin{equation*}
\blue{\lim_{t\to t_0} ||\u(t; \u_0)||_{H^m} = \infty, \qquad 0 < t_0 <
  \infty} \quad ?
\end{equation*}


\end{itemize}

\Emp
}

%%%%%%%%%%%%%%%%%%%%%%%%%%%%%%%%%%%%%%%%%%%%%%%%%%%%%%%%%%%%%%%%%%%%%%%%%%
%%%%%%%%%%%%%%%%%%%%%%%%%%%%%%%%%%%%%%%%%%%%%%%%%%%%%%%%%%%%%%%%%%%%%%%%%%
%%%%%%%%%%%%%%%%%%%%%%%%%%%%%%%%%%%%%%%%%%%%%%%%%%%%%%%%%%%%%%%%%%%%%%%%%%

\frame
{
%\frametitle{What is known? --- Available Estimates}
\hspace*{-1.0cm}
\Bmp{12cm}
\raggedright


\begin{itemize}
\item<1-> \blue{{\sc Beale-Kato-Majda (BKM) criterion}  (Beale et al.~1984)} --- 
A smooth solution $\u$ of the Euler system develops a singularity at
$t = t_0$ if and only if
\begin{equation*}
\blue{\lim_{t \rightarrow t_0}  \int_{0}^{t} ||\bomega(\tau)||_{L^\infty} \;d\tau = \infty,  \qquad \qquad \bomega = \bnabla\times \u}
\end{equation*}

\item<2-> The BKM criterion is related to a priori estimates based on Kato's theorem (Majda \& Bertozzi, 2002)
\begin{gather*}\hspace*{-0.3cm}
\visible<3->{
\int_{0}^{t} ||\bomega(\tau)||_{L^\infty} \;d\tau \leq t \sup_{0\leq \tau \leq t} ||\bnabla \u(\tau)||_{L^\infty}
\leq C\, t \sup_{0\leq \tau \leq t} ||\u(\tau)||_{H^m},\  m > \frac{5}{2}} \\
\visible<4->{
\blue{||\u(t)||_{H^m} \leq ||\u_0||_{H^m}\exp\left\{C_1 \exp\left(C_2\int_{0}^{t} ||\bomega(\tau)||_{L^\infty} \;d\tau \right)\right\}}}
\end{gather*}

\item<5->The double exponential on the RHS suggests a much more rapid growth of
$\| \u(t) \|_{H^m}$ than that of $\| \bomega(t) \|_{L^\infty}$, as $t \rightarrow t_0$. 


%\medskip
%\item<4-> {\blue{\sc Question:}} How to look for initial data $\u_0$
%  which may lead to singularity in finite time?

\end{itemize}

\Emp
}

%%%%%%%%%%%%%%%%%%%%%%%%%%%%%%%%%%%%%%%%%%%%%%%%%%%%%%%%%%%%%%%%%%%%%%%%%%
%%%%%%%%%%%%%%%%%%%%%%%%%%%%%%%%%%%%%%%%%%%%%%%%%%%%%%%%%%%%%%%%%%%%%%%%%%
%%%%%%%%%%%%%%%%%%%%%%%%%%%%%%%%%%%%%%%%%%%%%%%%%%%%%%%%%%%%%%%%%%%%%%%%%%

\frame
{
%\frametitle{What is known? --- Available Estimates}
\hspace*{-1.0cm}
\Bmp{12cm}
\raggedright


\begin{itemize}
\item<1->  {\blue{\sc Question:}} How to look for initial data $\u_0$ which may lead to singularity in finite time?
\medskip

\item<2-> The computational search for singularities in Euler flows
  has had a long history: Brachet et al.~(1983), Pumir \& Siggia
  (1990), Brachet (1991), Kerr (1993), Pelz (2001), Bustamante \& Kerr
  (2008), Matsumoto et al.~(2008), Ohkitani \& Constantin (2008),
  Ohkitani (2008), Grafke et al.~(2008), Gibbon et al.~(2008), Hou
  (2009), Siegel \& Caflisch (2009),Orlandi et al.~(2012), Bustamante
  \& Brachet (2012), Kerr (2013), Orlandi et al.~(2014), Campolina \&
  Mailybaev (2018), Larios et al.~(2018),\dots
\medskip

\item<3-> Well-documented blow-up in axisymmetric
Euler flows on a bounded cylindrical domain (Luo \& Hou, 2014)
\begin{itemize}
\item<4-> recent results by Hou (2022) also suggest singularity
formation away from the boundaries
\end{itemize}

\item<5->[$\Rightarrow$] \red{Search for singularities {\em systematically} using methods of PDE optimization}


\end{itemize}

\Emp
}

%%%%%%%%%%%%%%%%%%%%%%%%%%%%%%%%%%%%%%%%%%%%%%%%%%%%%%%%%%%%%%%%%%%%%%%%%%
%%%%%%%%%%%%%%%%%%%%%%%%%%%%%%%%%%%%%%%%%%%%%%%%%%%%%%%%%%%%%%%%%%%%%%%%%%
%%%%%%%%%%%%%%%%%%%%%%%%%%%%%%%%%%%%%%%%%%%%%%%%%%%%%%%%%%%%%%%%%%%%%%%%%%

%\section{Systematic Search for Singularities via PDE Optimization}
\subsection{Optimization Problems for Navier-Stokes}
\frame
{
\hspace*{-0.75cm}
\Bmp{12cm}
\raggedright

\begin{itemize}
\item<1-> Want to maximize enstrophy at time $T$, with $\E_0 :=
  \E(\u_0) > 0$ fixed, to see if\/\red{$\E_T(\u_0) := \E(\u(T;\u_0))$}\/can
  become infinite

\begin{problem}[1]
\vspace*{-0.5cm}
\blue{\small
\begin{align*}
& \max_{\u_0\in\mathcal{Q}_{\E_0}} \, \E_T(\u_0), \quad \text{where} \\
& \qquad \mathcal{Q}_{\E_0}  =  \left\{\u_0\in H^1(\Omega)\,\colon\,\bnabla\cdot\u_0 = 0, \; \int_{\Omega} \u_0 \, d\x = 0, \; \E(\u_0) = \E_0 \right\}, \\
& \text{subject to:} \quad  \left\{
\begin{aligned}
&\Dpartial{\u}{t}+({\u}\cdot{\bnabla}){\u}
+{\bnabla}{p}-\nu\Delta{{\u}}=\0,  && \text{in} \ \Omega\times(0,T] \\ 
&{\bnabla}\cdot{\u}=0,                     && \text{in} \ \Omega\times(0,T] \\ 
&\u = \u_0                 &&  \text{in} \ \Omega \ \text{at} \ t=0 \\
&\textrm{Periodic Boundary Conditions}                  &&  \text{on} \ \Gamma\times(0,T] 
\end{aligned}
%\label{eq:NS}
\right.
\end{align*}
}
\end{problem}
\medskip

\item<2->
A formidable, but solvable, PDE optimization problem


\end{itemize}

\Emp
}

%%%%%%%%%%%%%%%%%%%%%%%%%%%%%%%%%%%%%%%%%%%%%%%%%%%%%%%%%%%%%%%%%%%%%%%%%%
%%%%%%%%%%%%%%%%%%%%%%%%%%%%%%%%%%%%%%%%%%%%%%%%%%%%%%%%%%%%%%%%%%%%%%%%%%
%%%%%%%%%%%%%%%%%%%%%%%%%%%%%%%%%%%%%%%%%%%%%%%%%%%%%%%%%%%%%%%%%%%%%%%%%%
%\subsection{Solution Approach}
\frame
{
\hspace*{-1.0cm}
\Bmp{12cm}
\raggedright

\begin{itemize}
\item<1->
{\em Local} maximizers found via discretized gradient flow 
\footnotesize
\begin{equation*}
\begin{aligned}
\uET^{(n+1)} & =  \mathbb{P}_{\mathcal{Q}_{\E_0}}\left(\;\uET^{(n)} + \tau_n \blue{\nabla\E_T\left(\uET^{(n)}\right)}\;\right), \\ 
\uET^{(1)} & =  \u^0,
\end{aligned} \vspace*{-0.2cm}
\end{equation*}
where:
\medskip
\begin{itemize}
\small
\item<2-> \blue{$\nabla\E_T(\u_0)$}\/is the gradient of the objective
  functional $\E_T(\u_0)$ with respect to the initial data $\u_0$

\bigskip
\item<3->
\mbox{
  \Bmpt{7.0cm} \vspace*{-1.725cm}
step size $\tau^{(n)}$ is found via {\it arc minimization} and 
the projection on the constraint manifold $\mathcal{Q}_{\E_0}$ is given by
\begin{equation*}
\mathbb{P}_{\mathcal{Q}_{\E_0}}(\u_0) = \sqrt{\frac{\E_0}{\ET}}\,\u_0
\end{equation*}
\Emp \quad
\Bmp{5.0cm}
\vspace*{-0.5cm}
\visible<3->{\includegraphics[width=0.6\textwidth]{Figs9/normal_H1space}}
\Emp}

\end{itemize}

\end{itemize}

\Emp
}


%%%%%%%%%%%%%%%%%%%%%%%%%%%%%%%%%%%%%%%%%%%%%%%%%%%%%%%%%%%%%%%%%%%%%%%%%%
%%%%%%%%%%%%%%%%%%%%%%%%%%%%%%%%%%%%%%%%%%%%%%%%%%%%%%%%%%%%%%%%%%%%%%%%%%
%%%%%%%%%%%%%%%%%%%%%%%%%%%%%%%%%%%%%%%%%%%%%%%%%%%%%%%%%%%%%%%%%%%%%%%%%%
%\subsection{Sobolev Gradients}
\frame
{
\hspace*{-1.0cm}
\Bmp{12cm}
\raggedright

\begin{itemize}
\item<1-> How to ensure the required smoothness of the
  gradients\/\blue{$\nabla\E_T \in H^1$}?
  \smallskip
  
\item<2-> Defining the from {\it adjoint system}
 \footnotesize \blue{
\begin{align*}
 \L^*\begin{bmatrix} \u^* \\ p^* \end{bmatrix} := 
& \begin{bmatrix}
-\partial_{t}\u^*-\left[\bnabla\u^*+\bnabla\u^{*^T}\right]\u-\bnabla p^*-\nu\bDelta\u^* \\
-\bnabla\cdot\u^*
\end{bmatrix}  = \begin{bmatrix} \bDelta\u \\ 0\end{bmatrix}, \\
 \u^*(T)= & \mathbf{0} 
\end{align*}}
\small
the G\^{a}teaux differential of $\ET$ becomes \blue{$\E'_T(\u_0;\u_0') = \int_\Omega \u'_0\cdot\u^*(0)  \,d\x$}
\medskip

\item<3-> Since $\E_T'(\u_0,\cdot)$ is a bounded linear functional on
  $L^2(\Omega)$ and on $H^1(\Omega)$, the gradient can be
  deduced from the Riesz representation theorem
  \small
\begin{equation*}
\E'_T(\u_0;\u_0')
= \Big\langle \blue{\nabla^{L_2}\ET}, \u_0' \Big\rangle_{L^2(\Omega)}
= \Big\langle \red{\nabla\ET}, \u_0' \Big\rangle_{H^1(\Omega)},
%\label{eq:riesz}
\end{equation*}

\item<4-> Using the $L^2$ inner product: \qquad \blue{$\nabla^{L^2}\ET =
    \u^*(0)$}
  \smallskip
  
\item<5-> Using the $H^1$ inner product, an elliptic BVP is obtained:
  \small
  \begin{equation*}
 \red{ \left[ \Id \, - \,\ell_1^2 \,\bDelta \right] \nabla\E_T(\u_0)}
  = \blue{\nabla^{L_2} \E_T(\u_0)}   \qquad \text{in} \ \Omega
  \end{equation*}
  
 


\end{itemize}

\Emp
}


%%%%%%%%%%%%%%%%%%%%%%%%%%%%%%%%%%%%%%%%%%%%%%%%%%%%%%%%%%%%%%%%%%%%%%%%%%
%%%%%%%%%%%%%%%%%%%%%%%%%%%%%%%%%%%%%%%%%%%%%%%%%%%%%%%%%%%%%%%%%%%%%%%%%%
%%%%%%%%%%%%%%%%%%%%%%%%%%%%%%%%%%%%%%%%%%%%%%%%%%%%%%%%%%%%%%%%%%%%%%%%%%

\begin{frame}
\frametitle{Computational Algorithm}
%\textbf{Computational Algorithm:}
%\footnotesize
\begin{enumerate}
\setcounter{enumi}{-1}
\item[]<2->
\hspace*{-1.0cm} $\bullet$ set $\E_0$ and $T$ \\
\hspace*{-1.0cm} $\bullet$ provide initial guess for the initial data $\uET$ 
\medskip 
\item[1.]<3->
\blue{solve the Navier-Stokes system for $\{\u,p\}$}
\medskip 

\item[2.]<4->
\blue{solve the adjoint Navier-Stokes system for $\{\u^*,p^*\}$}
\medskip 

\item[3.]<5->
use $\u$ and $\u^*$ to compute $\bnabla^{L_2} \E_T$
\medskip 

\item[4.]<6->
determine the Sobolev gradient $\bnabla^{H^1} \E_T$
\medskip 

\item[5.]<7->
update the initial data while enforcing the enstrophy constraint 
\footnotesize
\begin{equation*}
\hspace*{2.5cm} \uET^{(n+1)} =  \mathbb{P}_{\mathcal{Q}_{\E_0}}\left(\;\uET^{(n)} + \tau_n \nabla\E_T\left(\uET^{(n)}\right)\;\right)
\end{equation*}
%\medskip 

\item[]<8->
\hspace*{-1.0cm} $\bullet$ iterate $1.$ through $5.$ until convergence, i.e. until 
\footnotesize
\begin{equation*}
\frac{\E(\uET^{(n+1)}) -  \E(\uET^{(n)})}{ \E(\uET^{(n)})} < \epsilon
\end{equation*}
\end{enumerate}
\normalsize
\end{frame}




%%%%%%%%%%%%%%%%%%%%%%%%%%%%%%%%%%%%%%%%%%%%%%%%%%%%%%%%%%%%%%%%%%%%%%%%%%
%%%%%%%%%%%%%%%%%%%%%%%%%%%%%%%%%%%%%%%%%%%%%%%%%%%%%%%%%%%%%%%%%%%%%%%%%%
%%%%%%%%%%%%%%%%%%%%%%%%%%%%%%%%%%%%%%%%%%%%%%%%%%%%%%%%%%%%%%%%%%%%%%%%%%

\frame
{
\vspace*{-0.1cm}
\hspace*{-1.1cm}
\Bmp{12cm}
\raggedright
\frametitle{Maximum Sustained Rate of  Enstrophy Growth $\frac{d\E}{dt} \sim C \, \E^\alpha$}
%\framesubtitle{$\Large \frac{d\E}{dt} \sim C \, \E^\alpha$}
\begin{figure}
%\linespread{1.1}
\setcounter{subfigure}{0}
\begin{center}
\includegraphics[width=0.75\textwidth]{Figs9/dEvsE_v1.eps}
\end{center}
\end{figure}
\footnotesize
\begin{center}
\Bmp{0.7\textwidth}
\hspace*{-0.1cm} \blue{-----} extreme trajectories with optimal initial data $\tuET$  \\ \bigskip
\red{-----} instantaneous maximizers  $\tuE$ 
\Emp
\end{center}
\Emp
}



%%%%%%%%%%%%%%%%%%%%%%%%%%%%%%%%%%%%%%%%%%%%%%%%%%%%%%%%%%%%%%%%%%%%%%%%%%
%%%%%%%%%%%%%%%%%%%%%%%%%%%%%%%%%%%%%%%%%%%%%%%%%%%%%%%%%%%%%%%%%%%%%%%%%%
%%%%%%%%%%%%%%%%%%%%%%%%%%%%%%%%%%%%%%%%%%%%%%%%%%%%%%%%%%%%%%%%%%%%%%%%%%
\frame
{
\vspace*{-0.1cm}
\hspace*{-1.1cm}
\Bmp{12cm}
\raggedright
\frametitle{Maximum enstrophy $\max_{\u_0} \mathcal{E}(T)$ versus $T$ for and $\E_0$}

\vspace*{-0.15cm}

\begin{figure}
\linespread{1.1}
\setcounter{subfigure}{0}
\begin{center}
\mbox{\hspace*{-0.15cm}
\includegraphics[width=0.575\textwidth]{Figs9/maxEt_vs_T_v4b.eps} 
%\quad
\hspace*{-0.8cm}
\visible<2->{
\includegraphics[width=0.575\textwidth]{Figs9/maxEt_vs_E0_v5b.eps}}}
\end{center}
\end{figure}

\visible<2->{
\bigskip
\centering \hspace*{0.75cm}
$\large \boxed{\blue{\max_T \max_{t \in [0,T]} \E(\tuET(t))  \ \sim \ \left( 0.224 \ \pm 0.006 \right) \, \E_0^{\red{{1.490 \, \pm 0.004}}} }}$}

\Emp
}


%%%%%%%%%%%%%%%%%%%%%%%%%%%%%%%%%%%%%%%%%%%%%%%%%%%%%%%%%%%%%%%%%%%%%%%%%%
%%%%%%%%%%%%%%%%%%%%%%%%%%%%%%%%%%%%%%%%%%%%%%%%%%%%%%%%%%%%%%%%%%%%%%%%%%
%%%%%%%%%%%%%%%%%%%%%%%%%%%%%%%%%%%%%%%%%%%%%%%%%%%%%%%%%%%%%%%%%%%%%%%%%%
%\subsection{Structure of the Flows and Vortex Reconnection}
\frame
{
\vspace*{-0.1cm}
\hspace*{-0.95cm}
\Bmp{12cm}
\raggedright
\frametitle{Structure of the optimal initial data $\tuET$}
\framesubtitle{($\E_0=500$, $T=0.017$)}
\vspace*{-0.2cm}

\begin{figure}
\setcounter{subfigure}{0}
\mbox{\subfigure[$\omega_x$]{\includegraphics[width=0.3\textwidth]{Figs9/E500T017_omega1.png}}\qquad
\subfigure[$\omega_y$]{\includegraphics[width=0.3\textwidth]{Figs9/E500T017_omega2.png}}\qquad
\subfigure[$\omega_z$]{\includegraphics[width=0.3\textwidth]{Figs9/E500T017_omega3.png}}}
\end{figure}

\Emp
}

%%%%%%%%%%%%%%%%%%%%%%%%%%%%%%%%%%%%%%%%%%%%%%%%%%%%%%%%%%%%%%%%%%%%%%%%%%
%%%%%%%%%%%%%%%%%%%%%%%%%%%%%%%%%%%%%%%%%%%%%%%%%%%%%%%%%%%%%%%%%%%%%%%%%%
%%%%%%%%%%%%%%%%%%%%%%%%%%%%%%%%%%%%%%%%%%%%%%%%%%%%%%%%%%%%%%%%%%%%%%%%%%

\frame
{
\vspace*{-0.1cm}
\hspace*{-0.95cm}
\Bmp{12cm}
\raggedright
\frametitle{Time evolution of the extremal flow}
\framesubtitle{($\E_0 = 500$ and $\tTE = 0.17$)}

\centering
\movie[width=0.75\textwidth,loop,showcontrols]{\includegraphics[width=0.75\textwidth]{Figs9/E500T017t0.png}}{../maxE3D/Submit/Movies/Movie1b.mov}
\Emp
}



%%%%%%%%%%%%%%%%%%%%%%%%%%%%%%%%%%%%%%%%%%%%%%%%%%%%%%%%%%%%%%%%%%%%%%%%%%
%%%%%%%%%%%%%%%%%%%%%%%%%%%%%%%%%%%%%%%%%%%%%%%%%%%%%%%%%%%%%%%%%%%%%%%%%%
%%%%%%%%%%%%%%%%%%%%%%%%%%%%%%%%%%%%%%%%%%%%%%%%%%%%%%%%%%%%%%%%%%%%%%%%%%

%\subsection{Optimization Problem Based on the LPS Condition}
\frame
{
\hspace*{-0.75cm}
\Bmp{12cm}
\raggedright

\small
\begin{itemize}
\item<1-> Maximize the  Ladyzhenskaya-Prodi-Serrin functional with $q=4$, $p=8$ \\
\vspace*{-0.25cm}
\begin{equation*} 
\red{\Phi_T(\u_0)  := \frac{1}{T} \int_0^T \| \u(\tau) \|_{L^4(\Omega)}^8 \, d\tau}
\end{equation*} \vspace*{-0.35cm}

\begin{problem}[2]
\vspace*{-0.75cm}
\blue{\small
\begin{align*}
& \max_{\u_0\in\mathcal{L}_{B}} \, \Phi_T(\u_0), \quad \text{where} \\
& \qquad \mathcal{L}_{B}  =  \left\{\u_0\in H^{3/4}(\Omega)\,\colon\,\bnabla\cdot\u_0 = 0, \; \int_{\Omega} \u_0 \, d\x = 0, \;  \|\u_0\|_{L^4(\Omega)} = B \right\}, \\
& \text{subject to:} \quad  \left\{
\begin{aligned}
&\Dpartial{\u}{t}+({\u}\cdot{\bnabla}){\u}
+{\bnabla}{p}-\nu\Delta{{\u}}=\0,  && \text{in} \ \Omega\times(0,T] \\ 
&{\bnabla}\cdot{\u}=0,                     && \text{in} \ \Omega\times(0,T] \\ 
&\u = \u_0                 &&  \text{in} \ \Omega \ \text{at} \ t=0 \\
&\textrm{Periodic Boundary Conditions}                  &&  \text{on} \ \Gamma\times(0,T] 
\end{aligned}
%\label{eq:NS}
\right.
\end{align*}
}
\end{problem}
%\medskip

\item<2-> Solutions sought in $H^{3/4}$, the largest Sobolev space with
  Hilbert structure embedded in $L^4$

\end{itemize}

\Emp
}

%%%%%%%%%%%%%%%%%%%%%%%%%%%%%%%%%%%%%%%%%%%%%%%%%%%%%%%%%%%%%%%%%%%%%%%%%%
%%%%%%%%%%%%%%%%%%%%%%%%%%%%%%%%%%%%%%%%%%%%%%%%%%%%%%%%%%%%%%%%%%%%%%%%%%
%%%%%%%%%%%%%%%%%%%%%%%%%%%%%%%%%%%%%%%%%%%%%%%%%%%%%%%%%%%%%%%%%%%%%%%%%%

\frame
{
\vspace*{-0.1cm}
\hspace*{-1.1cm}
\Bmp{12cm}
\raggedright
\frametitle{$\|\u(t)\|_{L^4}$ versus time $t$ \quad \hfill  \quad \hfill  $\max_{\u_0} \Phi_T(\u_0)$ versus $T$}

\vspace*{-0.15cm}

\begin{figure}
\linespread{1.1}
\setcounter{subfigure}{0}
\begin{center}
\mbox{\hspace*{-0.15cm}
\includegraphics[width=0.575\textwidth]{Figs9/LvsT_L12000.pdf} 
%\quad
\hspace*{-0.8cm}
\visible<2->{
\includegraphics[width=0.575\textwidth]{Figs9/LPS_vs_T_L.pdf}}}
\end{center}
\end{figure}

%\Bmp{0.9\textwidth}
\centering
\small
\blue{---$\blacksquare$---} partially symmetric branch \qquad\qquad\qquad 
\red{---\mbox{\Large $\bullet$}---} asymmetric branch \\
\bigskip
\normalsize
\centering
\red{No evidence for unbounded growth of $\|\u(t)\|_{L^4}$ and singularity formation}

%\Emp
\Emp
}

%%%%%%%%%%%%%%%%%%%%%%%%%%%%%%%%%%%%%%%%%%%%%%%%%%%%%%%%%%%%%%%%%%%%%%%%%%
%%%%%%%%%%%%%%%%%%%%%%%%%%%%%%%%%%%%%%%%%%%%%%%%%%%%%%%%%%%%%%%%%%%%%%%%%%
%%%%%%%%%%%%%%%%%%%%%%%%%%%%%%%%%%%%%%%%%%%%%%%%%%%%%%%%%%%%%%%%%%%%%%%%%%

\frame
{
\vspace*{-0.1cm}
\hspace*{-0.95cm}
\Bmp{12cm}
\raggedright
\frametitle{Time evolution of the extremal flow}
\framesubtitle{( $B^4=12{,}000$ and $T=0.01$)}

\centering
\movie[width=0.6\textwidth,loop,showcontrols]{\includegraphics[width=0.6\textwidth]{Figs9/Q12000T001_L_B4.png}}{../maxLPS/Submit/Movies/movie1.mov}
\Emp
}

%%%%%%%%%%%%%%%%%%%%%%%%%%%%%%%%%%%%%%%%%%%%%%%%%%%%%%%%%%%%%%%%%%%%%%%%%%
%%%%%%%%%%%%%%%%%%%%%%%%%%%%%%%%%%%%%%%%%%%%%%%%%%%%%%%%%%%%%%%%%%%%%%%%%%
%%%%%%%%%%%%%%%%%%%%%%%%%%%%%%%%%%%%%%%%%%%%%%%%%%%%%%%%%%%%%%%%%%%%%%%%%%

\frame
{
\vspace*{-0.1cm}
\hspace*{-0.95cm}
\Bmp{12cm}
\raggedright
\frametitle{Maximum enstrophy $\max_T \max_{\u_0} \mathcal{E}(T)$ vs. $\E_0$}

\vspace*{-0.5cm}
\centering
\includegraphics[width=0.7\textwidth]{Figs9/maxET_comp_v2.eps} 
\medskip

%\small
\red{- - -\mbox{\Large $\bullet$}- - -} solutions of Problem 1 \qquad\qquad
\blue{---\mbox{\Large $\diamond$}---} solutions of Problem 2 \\

\vspace*{-0.5cm}
\begin{equation*}
\boxed{\max_T \max_{\u_0} \mathcal{E}(T) \sim \E_0^{3/2}}
\end{equation*}

\Emp
}



%%%%%%%%%%%%%%%%%%%%%%%%%%%%%%%%%%%%%%%%%%%%%%%%%%%%%%%%%%%%%%%%%%%%%%%%%%
%%%%%%%%%%%%%%%%%%%%%%%%%%%%%%%%%%%%%%%%%%%%%%%%%%%%%%%%%%%%%%%%%%%%%%%%%%
%%%%%%%%%%%%%%%%%%%%%%%%%%%%%%%%%%%%%%%%%%%%%%%%%%%%%%%%%%%%%%%%%%%%%%%%%%

\subsection{Optimization Problem for Euler}
\frame
{
  \hspace*{-0.75cm}
\Bmp{12cm}
\raggedright

\begin{itemize}
\item<1-> The symmetry
  \blue{$\lambda \u_\lambda(\lambda t, \x) = \u(t, \x)$,
    $\lambda \neq 0$} implies  that \\
  the ``size'' of $\u_0$ and the time $T$ and not independent 
 \begin{itemize}
 \item<2->
   either $||\u_0||_{\dot{H}^3}$ or $T$ can be fixed
 \end{itemize}
 
  \medskip
  
\item<3-> Set $m = 3$ in Kato's theorem and maximize
  \begin{equation*}
  \red{\Phi_T(\u_0) = ||\u(T;\u_0)||_{\dot{H}^3} \qquad \text{subject to} \qquad  ||\u_0||_{\dot{H}^3} = 1}
  \end{equation*}
for different (increasing) $T$

\bigskip
\item<4->
  Unlike the Navier-Stokes system, the Euler system {\em does not}
  instantaneously regularize solutions (at $t = 0^+$) 
  \begin{itemize}
  \item<5-> to facilitate numerical solution of the Euler system, the
    optimal initial condition is sought in the Gevrey space $G^\sigma$
    of analytic functions endowed with the inner product \\ \vspace*{-0.2cm}
    \begin{equation*}
\blue{\langle \v, \u\rangle_{G^\sigma} = \sum_{\j \in \mathbb Z^3} 
(1+|2\pi\j|^2)^3 e^{4\pi\sigma|\j|} \hat{\v}_{\j} \cdot
\overline{\hat{\u}}_{\j}, \qquad \sigma > 0}
\end{equation*}
  \end{itemize}
\end{itemize}

\Emp
}

%%%%%%%%%%%%%%%%%%%%%%%%%%%%%%%%%%%%%%%%%%%%%%%%%%%%%%%%%%%%%%%%%%%%%%%%%%
%%%%%%%%%%%%%%%%%%%%%%%%%%%%%%%%%%%%%%%%%%%%%%%%%%%%%%%%%%%%%%%%%%%%%%%%%%
%%%%%%%%%%%%%%%%%%%%%%%%%%%%%%%%%%%%%%%%%%%%%%%%%%%%%%%%%%%%%%%%%%%%%%%%%%

\frame
{
\hspace*{-0.75cm}
\Bmp{12cm}
\raggedright

\begin{itemize}
\item<1->[]
\begin{problem}[3]
\vspace*{-0.5cm}
\blue{\small
\begin{align*}
& \max_{\u_0\in\mathcal{M}_{1}} \, \Phi_T(\u_0), \quad \text{where} \\
& \qquad \mathcal{M}_{1}  =  \left\{\u_0\in G^{\sigma}\,\colon\,\bnabla\cdot\u_0 = 0, \; \int_{\Omega} \u_0 \, d\x = 0, \; ||\u_0||_{\dot{H}^3} = 1 \right\}, \\
& \text{subject to:} \quad  \left\{
\begin{aligned}
&\Dpartial{\u}{t}+({\u}\cdot{\bnabla}){\u}
+{\bnabla}{p}=\0,  && \text{in} \ \Omega\times(0,T] \\ 
&{\bnabla}\cdot{\u}=0,                     && \text{in} \ \Omega\times(0,T] \\ 
&\u = \u_0                 &&  \text{in} \ \Omega \ \text{at} \ t=0 \\
&\textrm{Periodic Boundary Conditions}                  &&  \text{on} \ \Gamma\times(0,T] 
\end{aligned}
%\label{eq:NS}
\right.
\end{align*}
}
\end{problem}

\item<2->
Solution of the forward Euler and the backward adjoint system
facilitated by time-reversibility


\end{itemize}

\Emp
}

%%%%%%%%%%%%%%%%%%%%%%%%%%%%%%%%%%%%%%%%%%%%%%%%%%%%%%%%%%%%%%%%%%%%%%%%%%
%%%%%%%%%%%%%%%%%%%%%%%%%%%%%%%%%%%%%%%%%%%%%%%%%%%%%%%%%%%%%%%%%%%%%%%%%%
%%%%%%%%%%%%%%%%%%%%%%%%%%%%%%%%%%%%%%%%%%%%%%%%%%%%%%%%%%%%%%%%%%%%%%%%%%

\frame
{
\hspace*{-0.75cm}
\Bmp{12cm}
\raggedright


  \blue{$\teta^N_T$ --- optimal initial condition obtained at time $T$ with
  resolution $N$}

\bigskip
\begin{itemize}
\item<2-> ``short'' times $T$:  \qquad $\lim_{N\to\infty} \Phi_{T}\left(\teta_T^N\right) < \infty$ 
\item<2->[]
$\Longrightarrow$ Euler system is well-posed on $[0, T]$
\bigskip
\item<3->  ``long'' times $T$: \qquad $\lim_{N\to\infty} \Phi_{T}\left(\teta_T^N\right) = \infty$
\item<3->[]
 $\Longrightarrow$ \red{a singularity possible at $t_0 \leq T$}
\bigskip\bigskip
\item<4-> In practice, we use $N = 128, 256, 512, 1024$

\end{itemize}

\Emp
}


%%%%%%%%%%%%%%%%%%%%%%%%%%%%%%%%%%%%%%%%%%%%%%%%%%%%%%%%%%%%%%%%%%%%%%%%%%
%%%%%%%%%%%%%%%%%%%%%%%%%%%%%%%%%%%%%%%%%%%%%%%%%%%%%%%%%%%%%%%%%%%%%%%%%%
%%%%%%%%%%%%%%%%%%%%%%%%%%%%%%%%%%%%%%%%%%%%%%%%%%%%%%%%%%%%%%%%%%%%%%%%%%

\frame
{
\hspace*{-0.75cm}
\Bmp{12cm}
\raggedright

\begin{figure}
  \centering
  \mbox{
  \subfigure[$T = 25$]{\includegraphics[width=0.5\textwidth]{../Euler_Singularity/Figs/time_25_taylor_resol_obj_compare}}
  \subfigure[$T = 75$]{\includegraphics[width=0.5\textwidth]{../Euler_Singularity/Figs/time_75_optimization_resol_obj_compare}}}     
     \end{figure}

     \centering \small
     Objective functional $\Phi_{T}\left(\teta_T^N\right)$ obtained
     for different time windows $T$ with different spatial resolutions
     $N$

\Emp
}


%%%%%%%%%%%%%%%%%%%%%%%%%%%%%%%%%%%%%%%%%%%%%%%%%%%%%%%%%%%%%%%%%%%%%%%%%%
%%%%%%%%%%%%%%%%%%%%%%%%%%%%%%%%%%%%%%%%%%%%%%%%%%%%%%%%%%%%%%%%%%%%%%%%%%
%%%%%%%%%%%%%%%%%%%%%%%%%%%%%%%%%%%%%%%%%%%%%%%%%%%%%%%%%%%%%%%%%%%%%%%%%%

\frame
{
\hspace*{-0.75cm}
\Bmp{12cm}
\raggedright
\setcounter{subfigure}{0}

\begin{figure}
  \centering
  \mbox{
  \includegraphics[width=0.5\textwidth]{../Euler_Singularity/Figs/time_75_df_compare}
  \includegraphics[width=0.5\textwidth]{../Euler_Singularity/Figs/time_75_alpha_compare}}     
\end{figure}
     
     \centering \small
     The growth rate of $\|\u(t)\|_{\dot H^3}$:
     \smallskip
     \begin{equation*}
     \blue{ \frac{d \|\u(t)\|_{\dot H^3}}{dt} = C \|\u(t)\|_{\dot
       H^3}^{\red{\alpha}} \quad \Longrightarrow \quad
 \ln \left(\frac{d \|\u(t)\|_{\dot H^3}}{dt}\right) = \ln (C) +
 \red{\alpha} \ln \left(\|\u(t)\|_{\dot H^3}\right) }
\end{equation*}
\smallskip

\red{Computations become under-resolved at $t \approx  51.25$}

\Emp
}


%%%%%%%%%%%%%%%%%%%%%%%%%%%%%%%%%%%%%%%%%%%%%%%%%%%%%%%%%%%%%%%%%%%%%%%%%%
%%%%%%%%%%%%%%%%%%%%%%%%%%%%%%%%%%%%%%%%%%%%%%%%%%%%%%%%%%%%%%%%%%%%%%%%%%
%%%%%%%%%%%%%%%%%%%%%%%%%%%%%%%%%%%%%%%%%%%%%%%%%%%%%%%%%%%%%%%%%%%%%%%%%%

\frame
{
\hspace*{-0.75cm}
\Bmp{12cm}
\raggedright
\setcounter{subfigure}{0}

\begin{figure}
  \centering
  \mbox{
    \subfigure[$\omega_x$]{\includegraphics[width=0.45\textwidth]{../Euler_Singularity/misc/vort_x_0_shift}} \qquad
    \subfigure[$\omega_y$]{\includegraphics[width=0.45\textwidth]{../Euler_Singularity/misc/vort_y_0_shift}}}
  \subfigure[$\omega_z$]{\includegraphics[width=0.45\textwidth]{../Euler_Singularity/misc/vort_z_0_shift}}
     \end{figure}

     \centering \small \vspace*{-0.3cm}
     Vorticity components of the optimal initial condition
     $\teta_{75}^{1024}$
     
\Emp
}


%%%%%%%%%%%%%%%%%%%%%%%%%%%%%%%%%%%%%%%%%%%%%%%%%%%%%%%%%%%%%%%%%%%%%%%%%%
%%%%%%%%%%%%%%%%%%%%%%%%%%%%%%%%%%%%%%%%%%%%%%%%%%%%%%%%%%%%%%%%%%%%%%%%%%
%%%%%%%%%%%%%%%%%%%%%%%%%%%%%%%%%%%%%%%%%%%%%%%%%%%%%%%%%%%%%%%%%%%%%%%%%%

\frame
{
\hspace*{-0.75cm}
\Bmp{12cm}
\raggedright
\setcounter{subfigure}{0}

\begin{figure}
  \centering
  \mbox{\hspace*{-0.35cm}
     \movie[width=0.5\textwidth,loop,showcontrols]{\includegraphics[width=0.5\textwidth]{../Euler_Singularity/Figs/vort_mag_75}}{../Euler_Singularity/Submit/Movies/movie1.mov}
     \quad
     \movie[width=0.5\textwidth,loop,showcontrols]{\includegraphics[width=0.5\textwidth]{../Euler_Singularity/Figs/vort_proj_75}}{../Euler_Singularity/Submit/Movies/movie2.mov}}
%  \includegraphics[width=0.48\textwidth]{../Euler_Singularity/Figs/vort_mag_75} \quad
%  \includegraphics[width=0.48\textwidth]{../Euler_Singularity/Figs/vort_proj_75}}     
\end{figure}

\bigskip\bigskip
     \centering \small
     The vorticity field $\bnabla \times \u\left(75;
       \teta_{75}^{1024}\right)$ in the Euler flow corresponding to
     the optimal initial condition $ \teta_{75}^{1024}$

\Emp
}

%%%%%%%%%%%%%%%%%%%%%%%%%%%%%%%%%%%%%%%%%%%%%%%%%%%%%%%%%%%%%%%%%%%%%%%%%%
%%%%%%%%%%%%%%%%%%%%%%%%%%%%%%%%%%%%%%%%%%%%%%%%%%%%%%%%%%%%%%%%%%%%%%%%%%
%%%%%%%%%%%%%%%%%%%%%%%%%%%%%%%%%%%%%%%%%%%%%%%%%%%%%%%%%%%%%%%%%%%%%%%%%%
\frame
{
\frametitle{Conclusions}
\vspace{-0.1cm}
%\vspace*{0.0cm}
\hspace*{-1.0cm}
\Bmp{12cm}
\raggedright
\small
\begin{itemize}
\item<1-> In the extreme Navier-Stokes flows the enstrophy $\E(t)$ and the norm
  $\|\u(t)\|_{L^4}$ remain finite at all times
\begin{itemize}
\item<2-> \red{hence, even in such worst-case scenario there is no evidence
  for formation of singularity in finite time}

\smallskip
\item<3-> 
\blue{however, we do not know if the maximizers found are global}


\smallskip
\item<4-> the extreme behavior in the two cases is realized by
  entirely different mechanisms

\smallskip
\item<5-> the scaling of the maximum growth of enstrophy with $\E_0$
  is the same in both cases and, remarkably, the same as in 1D Burgers
  flows

\end{itemize}

\medskip
\item<6-> The extreme Euler flows
  \begin{itemize}
\item<7-> remain regular during sufficiently short times ($0 \le t  \le 25$)

  \smallskip
\item<8-> \blue{appear to form a singularity at a later time ($t >
    51.25$)} \\ \smallskip
  \visible<9->{\red{$\Longrightarrow$ the singular flow structure is nearly axisymmetric!}}
  
\end{itemize}
  
\end{itemize}
\Emp
}

%%%%%%%%%%%%%%%%%%%%%%%%%%%%%%%%%%%%%%%%%%%%%%%%%%%%%%%%%%%%%%%%%%%%%%%%%%
%%%%%%%%%%%%%%%%%%%%%%%%%%%%%%%%%%%%%%%%%%%%%%%%%%%%%%%%%%%%%%%%%%%%%%%%%%
%%%%%%%%%%%%%%%%%%%%%%%%%%%%%%%%%%%%%%%%%%%%%%%%%%%%%%%%%%%%%%%%%%%%%%%%%%
%\subsection{Relevant Energy-type Estimates}
\frame
{
\frametitle{Summary of Relevant Energy-type Estimates}
\vspace*{-0.5cm}
%\begin{itemize}
%\item<1->
\begin{center}
\hspace*{-1.1cm}
\begin{tabular}{l|c|c}      
  &  \Bmp{3.5cm} \small \begin{center} {\sc Best Estimate} \\ \smallskip \end{center} \Emp   
  & \Bmp{1.5cm} \small \begin{center} {\sc Sharp?}  \end{center} \Emp \\  \hline
 \Bmp{2.5cm}  \footnotesize {\begin{center} \smallskip 1D Burgers  \\ instantaneous \smallskip \end{center}} \Emp &  
  \footnotesize {$\frac{d\E(t)}{dt} \leq \frac{3}{2}\left(\frac{1}{\pi^2\nu}\right)^{1/3}\E(t)^{5/3}$}  & \Bmp{3.0cm} \footnotesize \red{\begin{center} \smallskip {\sc Yes} \\ Lu \& Doering (2008) \smallskip  \end{center}} \Emp \\ \hline 
  \Bmp{2.5cm} \footnotesize {\begin{center} \smallskip 1D Burgers  \\ finite--time \smallskip \end{center}} \Emp &  
   \footnotesize {$\max_{t \in [0,T]} \E(t) \leq C\,\E_0^{3/2}$} &
                                                                   \Bmp{3.5cm}
                                                                   \footnotesize
                                                                   \red{\begin{center}
                                                                       \smallskip
                                                                       {\sc
                                                                         Yes} \                                                                      
                                                                       Albritton
                                                                       \&
                                                                       De
                                                                       Nitti
                                                                       (2023)
                                                                       \smallskip  \end{center}} \Emp \\ \hline 
  \Bmp{2.5cm} \footnotesize \begin{center} \smallskip 2D Navier--Stokes  \\ instantaneous \smallskip\end{center} \Emp &  \Bmp{5.0cm} \footnotesize \centering\vspace{0.1cm} 
  $\frac{d\P(t)}{dt} \le -\left(\frac{\nu}{\E}\right)\P^2 + C_1\left(\frac{\E}{\nu}\right)\P$ \\ \vspace{0.1cm}
  $\frac{d\P(t)}{dt} \le \frac{C_2}{\nu} \K^{1/2}\, \P^{3/2}$ \\ \vspace{0.1cm} \Emp & \Bmp{3.0cm} \footnotesize \red{\begin{center} \smallskip {\sc Yes}  \\ Ayala \& P. (2013) \\ Ayala, Doering \& Simon (2017) \smallskip  \end{center}} \Emp  \\ \hline 
  \Bmp{2.5cm} \footnotesize \begin{center} \smallskip 2D Navier--Stokes \\ finite--time \smallskip\end{center} \Emp &  \footnotesize
   $\max_{t > 0} \P(t) \, \le \, \left[\P^{1/2}_0 + \frac{C_2}{4\nu^2} \K^{1/2}_0 \E_0\right]^2$ & \Bmp{3.0cm} \footnotesize \red{\begin{center} \smallskip {\sc Yes}  \\ Ayala \& P. (2013) \smallskip  \end{center}} \Emp   \\ \hline 
  \Bmp{2.5cm} \footnotesize {\begin{center} \smallskip 3D Navier--Stokes  \\ instantaneous  \smallskip \end{center}} \Emp &  
  \footnotesize {$\frac{d\E(t)}{dt} \le \frac{27 C^2}{128 \nu^3} \E(t)^3$} & \Bmp{3.0cm} \footnotesize \red{\begin{center} \smallskip {\sc Yes} \\ Lu \& Doering (2008) \smallskip  \end{center}} \Emp  \\ \hline 
  \Bmp{2.5cm} \footnotesize {\begin{center} \smallskip 3D Navier--Stokes  \\ finite--time \smallskip \end{center}} \Emp &  \Bmp{3.0cm} \footnotesize \vspace*{-0.2cm} {$\E(t) \le \frac{\E(0)}{\sqrt{1 - 4 \frac{C \E(0)^2}{\nu^3} t}}$} \Emp  &  
  \Bmp{2.0cm} \centering \footnotesize \smallskip \red{{\sc No} (???) \\ Kang, Yun \& P, (2020)} \smallskip\smallskip   \Emp
\end{tabular}
\end{center}
}

%%%%%%%%%%%%%%%%%%%%%%%%%%%%%%%%%%%%%%%%%%%%%%%%%%%%%%%%%%%%%%%%%%%%%%%%%%
%%%%%%%%%%%%%%%%%%%%%%%%%%%%%%%%%%%%%%%%%%%%%%%%%%%%%%%%%%%%%%%%%%%%%%%%%%
%%%%%%%%%%%%%%%%%%%%%%%%%%%%%%%%%%%%%%%%%%%%%%%%%%%%%%%%%%%%%%%%%%%%%%%%%%


%%\section[Agenda]{}
%\frame{
%%\frametitle{Review Article}

%\centering
%\includegraphics[width=1.0\textwidth]{Figs9/p21a_v2.png} \\
%\bigskip\medskip
%\small 
%Special Issue of \\  {\em Philosophical Transactions of the Royal Society A} \\
%{\bf{``Mathematical Problems in Physical Fluid Dynamics''}} \\
%Eds.~C.~R.~Doering, D.~Goluskin, B.~Protas \& J.-L.~Thiffeault 


%}



%%%%%%%%%%%%%%%%%%%%%%%%%%%%%%%%%%%%%%%%%%%%%%%%%%%%%%%%%%%%%%%%%%%%%%%%%%
%%%%%%%%%%%%%%%%%%%%%%%%%%%%%%%%%%%%%%%%%%%%%%%%%%%%%%%%%%%%%%%%%%%%%%%%%%
%%%%%%%%%%%%%%%%%%%%%%%%%%%%%%%%%%%%%%%%%%%%%%%%%%%%%%%%%%%%%%%%%%%%%%%%%%
\section{Inviscid instabilities of the 2D Taylor-Green vortex}
\subsection{Spectrum of the Linearized Euler Operator}
\frame{
\centering

\LARGE
{\sc PART II: \\ \medskip Inviscid instabilities of \\ the 2D Taylor-Green vortex}

\normalsize
\bigskip\bigskip
Joint work with:
\begin{itemize}
\item  Xinyu Zhao (post-doc, Assistant Professor at the New Jersey Institute of Technology as of September 2024)
\medskip
\item Roman Shvydkoy (University of Illinois at Chicago)
\end{itemize}
}


%%%%%%%%%%%%%%%%%%%%%%%%%%%%%%%%%%%%%%%%%%%%%%%%%%%%%%%%%%%%%%%%%%%%%%%%%%
%%%%%%%%%%%%%%%%%%%%%%%%%%%%%%%%%%%%%%%%%%%%%%%%%%%%%%%%%%%%%%%%%%%%%%%%%%
%%%%%%%%%%%%%%%%%%%%%%%%%%%%%%%%%%%%%%%%%%%%%%%%%%%%%%%%%%%%%%%%%%%%%%%%%%

\frame{
%\frametitle{References}
\centering
\medskip
\includegraphics[width=1.0\textwidth]{Figs9/zps24a2.png}
\bigskip\smallskip
}


%%%%%%%%%%%%%%%%%%%%%%%%%%%%%%%%%%%%%%%%%%%%%%%%%%%%%%%%%%%%%%%%%%%%%%%%%%
%%%%%%%%%%%%%%%%%%%%%%%%%%%%%%%%%%%%%%%%%%%%%%%%%%%%%%%%%%%%%%%%%%%%%%%%%%
%%%%%%%%%%%%%%%%%%%%%%%%%%%%%%%%%%%%%%%%%%%%%%%%%%%%%%%%%%%%%%%%%%%%%%%%%%

\frame
{
\hspace*{-0.75cm}
\Bmp{12cm}
\raggedright

\begin{itemize}
\item<1-> 2D Euler system on a periodic domain
\begin{equation*}
\begin{aligned}
\partial_t \omega + (\bds u \cdot \nabla)\omega = 0&, &\qquad &(\bds x, t) \in 
 \mbb T^2 \times (0, T], \\
 \omega|_{t = 0}  = \omega_0 &, &\qquad &\bds x \in \mbb T^2,
\end{aligned}
\end{equation*}
where 
\begin{itemize}
\item<2-> $\blue{\omega  = \nabla\times \bds u}$ is the scalar vorticity field satisfying 
$\int_{\mbb T^2} \omega \; d\bds x = 0$
\smallskip
\item<2-> Biot-Savart law: $\bds u = \nabla^\perp \Delta^{-1} \omega, \qquad \nabla^\perp = (-\partial_y, \partial_x)$
\end{itemize}

\smallskip
\item<3-> An equilibrium solution --- the 2D Taylor-Green vortex: 

\mbox{
\Bmpb{0.5\textwidth}
%\bigskip
\vspace*{-1.5cm}
\blue{
\begin{align*} 
\bds u_s & = (-\cos(x_1)\sin(x_2), \sin(x_1)\cos(x_2)) \\ \\
\omega_s & = 2\cos(x_1)\cos(x_2)
\end{align*}}
\Emp
\Bmpb{0.5\textwidth}
\centering
\visible<3->{\includegraphics[width=0.8\textwidth] {./Figs9/TaylorGreen.png}}
\Emp} 

\end{itemize}

     
\Emp
}

%%%%%%%%%%%%%%%%%%%%%%%%%%%%%%%%%%%%%%%%%%%%%%%%%%%%%%%%%%%%%%%%%%%%%%%%%%
%%%%%%%%%%%%%%%%%%%%%%%%%%%%%%%%%%%%%%%%%%%%%%%%%%%%%%%%%%%%%%%%%%%%%%%%%%
%%%%%%%%%%%%%%%%%%%%%%%%%%%%%%%%%%%%%%%%%%%%%%%%%%%%%%%%%%%%%%%%%%%%%%%%%%
\frame
{
\hspace*{-0.75cm}
\Bmp{12cm}
\raggedright

\begin{itemize}
\item<1-> Linearize the 2D Euler system around a steady solution $(\omega_s, \bds u_s)$ $\Longrightarrow$
\begin{align*}
\partial_t w &= {\blue \L} w, \qquad w|_{t = 0} = w_0, \qquad \text{where} \\
{\color {blue} \L w} :&= -(\bds u_s\cdot \nabla) w -(\bds u \cdot \nabla) \omega_s \\
& = {\color{blue} \left[-\bds u_s \cdot \nabla - \nabla\omega_s\cdot (\nabla^\perp \Delta^{-2})\right] w}, \qquad w \in H^m_0(\mbb T^2)
\end{align*}

%\smallskip
\item<2-> Linear stability of the equilibria determined by the spectrum of the operator $\L$
\blue{
\begin{align*}
\sigma(\L) := & \left\{ \lambda \in \CC \; : \; \text{$(\L - \lambda)$ is not bounded from below in $\mathcal{X}$} \right\} \\ 
= & \ \sigma_{\text{disc}}(\L) \cup \sigma_{\text{ess}} (\L)
\end{align*}}
\small \vspace*{-0.4cm}
\begin{itemize}
\item<3-> A point $z  \in \sigdis(\L)$ if
\begin{itemize}
\item<3->[(1)] $z$ is an isolated point in $\sigma(\L)$;
\item<3->[(2)]$z$ has finite multiplicity, i.e., $\bigcup_{r=1}^\infty
    \mathrm{Ker}(z - \L)^r$ is finite dimensional;
\item<3->[(3)] The range of $z-\L$ is closed. 
\end{itemize}
\smallskip
\item<4-> Otherwise, $z  \in \sigess(\L)$
\end{itemize}

\end{itemize}

     
\Emp
}

%%%%%%%%%%%%%%%%%%%%%%%%%%%%%%%%%%%%%%%%%%%%%%%%%%%%%%%%%%%%%%%%%%%%%%%%%%
%%%%%%%%%%%%%%%%%%%%%%%%%%%%%%%%%%%%%%%%%%%%%%%%%%%%%%%%%%%%%%%%%%%%%%%%%%
%%%%%%%%%%%%%%%%%%%%%%%%%%%%%%%%%%%%%%%%%%%%%%%%%%%%%%%%%%%%%%%%%%%%%%%%%%
\frame
{
\hspace*{-0.75cm}
\Bmp{12cm}
\raggedright

\begin{itemize}
\item<1-> 
The set of eigenvalues of $\L$ forms its point spectrum
\begin{equation*}
\blue{{\sigp(\L) := \left\{ \lambda \in \CC \; : \; \exists \phi (\x) \neq 0,  \quad \L \phi (\x) = \lambda \phi (\x), \quad \x \in \TT^2  \right\}}}
\end{equation*}

\bigskip
\item<2->
In finite dimensions linear operators can only have point spectrum

\bigskip
\item<3-> In infinite dimensions the situation is complicated by the
  presence of the essential spectrum $\sigess(\L)$, which may also contain
  eigenvalues of $\L$ that are not in $\sigdis(\L)$

\end{itemize}

     
\Emp
}


%%%%%%%%%%%%%%%%%%%%%%%%%%%%%%%%%%%%%%%%%%%%%%%%%%%%%%%%%%%%%%%%%%%%%%%%%%
%%%%%%%%%%%%%%%%%%%%%%%%%%%%%%%%%%%%%%%%%%%%%%%%%%%%%%%%%%%%%%%%%%%%%%%%%%
%%%%%%%%%%%%%%%%%%%%%%%%%%%%%%%%%%%%%%%%%%%%%%%%%%%%%%%%%%%%%%%%%%%%%%%%%%

\frame
{
\hspace*{-0.75cm}
\Bmp{12cm}
\raggedright

\begin{itemize}
\item<1-> For the 2D linear Euler operator $\L$ the essential spectrum
  $\sigess(\L)$ is well understood (Friedlander et al.~1997, Shvydkoy
  \& Latushkin 2003)
\begin{equation*}
  \blue{\sigess(\L; H^{m}_0) = \{ z\in \CC: |\Re(z)| \leq |m| \, \mumax \}},
\end{equation*}
where 
\begin{equation*}
\mumax :=  \lim_{t\to\infty} \frac{1}{t}\log \sup_{\bds x\in\mbb T^2} ||\nabla \varphi_t(\bds x)||
\end{equation*}
is the maximal Lyapunov exponent corresponding
to the Lagrangian flow $\varphi_t: \bxi \to \bds x(t; \bxi)$ generated
by the steady state via $\partial_t \bds x(t) = \bds u_s (\bds x(t))$

\medskip
\item<2-> If $z \in \sigess(\L)$, then there exists a sequence of
  weak-null unit vectors
  ${\{f_n\} \in H^m_0(\TT^2)}$,\/\blue{approximate eigenvectors,}\/such
  that
  \blue{$$\|(\L - z) f_n \|_{H^m} \rightarrow 0, \qquad n \rightarrow
    \infty$$} and $\{f_n\}$ does not contain any convergent
  subsequence
\end{itemize}

     
\Emp
}

%%%%%%%%%%%%%%%%%%%%%%%%%%%%%%%%%%%%%%%%%%%%%%%%%%%%%%%%%%%%%%%%%%%%%%%%%%
%%%%%%%%%%%%%%%%%%%%%%%%%%%%%%%%%%%%%%%%%%%%%%%%%%%%%%%%%%%%%%%%%%%%%%%%%%
%%%%%%%%%%%%%%%%%%%%%%%%%%%%%%%%%%%%%%%%%%%%%%%%%%%%%%%%%%%%%%%%%%%%%%%%%%

\frame
{
\hspace*{-0.75cm}
\Bmp{12cm}
\raggedright

\begin{itemize}
\item<1->
The the maximal Lyapunov exponent is determined by the hyperbolic stagnation point 
\begin{equation*}
\mumax = \max\left[ \lambda\left(\bnabla\u_s({\x_s}) \right)\right]
\end{equation*}

\medskip
\item<2->
  In the Taylor-Green vortex: \qquad \red{$\mumax = 1$} \\
  \medskip
  
  \centering
\visible<2->{\includegraphics[width=0.6\textwidth] {./Figs9/TaylorGreen.png}}

\end{itemize}

     
\Emp
}



%%%%%%%%%%%%%%%%%%%%%%%%%%%%%%%%%%%%%%%%%%%%%%%%%%%%%%%%%%%%%%%%%%%%%%%%%%
%%%%%%%%%%%%%%%%%%%%%%%%%%%%%%%%%%%%%%%%%%%%%%%%%%%%%%%%%%%%%%%%%%%%%%%%%%
%%%%%%%%%%%%%%%%%%%%%%%%%%%%%%%%%%%%%%%%%%%%%%%%%%%%%%%%%%%%%%%%%%%%%%%%%%

\frame
{
\hspace*{-0.75cm}
\Bmp{12cm}
\raggedright

\begin{itemize}
\item<1-> No general results available about the discrete spectrum $\sigdis(\L)$

\smallskip
\item<2-> Unstable eigenvalues found for shear flows (Drazin \& Reid
  1981; Friedlander et al.~1997), Kolmogorov flows (Friedlander et
  al.~1997) and the cellular ``cat’s-eye'' flow (Friedlander et
  al.~2000).

\medskip
\item<3->
  The problem remains open for the Taylor-Green vortex
  \smallskip
\begin{itemize}
\item<4-> Some work in the viscous setting (Sipp \&
  Jacquin 1998; Leblanc \& Godeferd 1999; Hattori \& Hirota
  2023). However, in that case the vortex decays as
  $\mathcal{O}(e^{-\nu t})$ and\/\blue{$\sigess(\L) = \emptyset$}
\end{itemize}

\smallskip
\item<5->
{\red{\sc Question:}} \ \blue{does $\L$ have unstable eigenvalues and, if
  so, do they belong to $\sigdis (\L)$ or $\sigess (\L)$?} \\ \vspace*{-0.3cm}
  \Bmp{\textwidth}
  \hspace*{5.0cm} 
\visible<5->{\includegraphics[width=0.35\textwidth]{Figs12/piess_01.pdf}}
\Emp
\end{itemize}

     
\Emp
}




%%%%%%%%%%%%%%%%%%%%%%%%%%%%%%%%%%%%%%%%%%%%%%%%%%%%%%%%%%%%%%%%%%%%%%%%%%
%%%%%%%%%%%%%%%%%%%%%%%%%%%%%%%%%%%%%%%%%%%%%%%%%%%%%%%%%%%%%%%%%%%%%%%%%%
%%%%%%%%%%%%%%%%%%%%%%%%%%%%%%%%%%%%%%%%%%%%%%%%%%%%%%%%%%%%%%%%%%%%%%%%%%

\frame
{
\hspace*{-0.75cm}
\Bmp{12cm}
\raggedright
\footnotesize
\begin{theorem}[Lin 2004]
  Suppose there exists an exponentially growing solution
  $e^{\lambda t} w_0$ of the linearized system $\partial_t w = \L w$
  with $\Re(\lambda) > 0$ and let\/\blue{$w_0 \in L^2(\TT^2)$.} Then we have
  the following
\begin{itemize}
\item[(i)] [regularity of growing modes] \red{$w_0 \in W^{1,p}(\TT^2)
  \bigcap L^q(\TT^2)$} for all $1 \le p < p^*$ and $1 \le q < \infty$,
  where
  \begin{equation*}
    p^* = \begin{cases}
      \red{\frac{\mumax}{\mumax - \Re(\lambda)}} = \frac{1}{1 -\Re(\lambda)}, \qquad & \mumax > \Re(\lambda) \\
      \infty, & \mumax \le \Re(\lambda)
      \end{cases}
    \end{equation*}

  \item[(ii)] [nonlinear instability] if there exists $\epsilon > 0$,
    such that for any $\delta > 0$ there is a solution $w^\delta(t)$
    of the 2D Euler system corresponding to the initial condition
    $\omega_0^{\delta}$, such that
\begin{equation*}
 \| \omega_0^{\delta} - \omega_s \|_{L^q} +  \| \omega_0^{\delta} - \omega_s \|_{W^{1,p}} \le \delta,
\end{equation*}
and
\begin{equation*}
{\sup_{0 < t < T_\delta} \|\omega^\delta(t) - \omega_s\|_{H^m} \geq \epsilon,}
\end{equation*}
%
where $T_\delta = O(|\ln(\delta)|)$, $q \in [1, \infty)$, $p =
[1,p^*)$ and $m \ge -1$.
\end{itemize}
\end{theorem}

     
\Emp
}

%%%%%%%%%%%%%%%%%%%%%%%%%%%%%%%%%%%%%%%%%%%%%%%%%%%%%%%%%%%%%%%%%%%%%%%%%%
%%%%%%%%%%%%%%%%%%%%%%%%%%%%%%%%%%%%%%%%%%%%%%%%%%%%%%%%%%%%%%%%%%%%%%%%%%
%%%%%%%%%%%%%%%%%%%%%%%%%%%%%%%%%%%%%%%%%%%%%%%%%%%%%%%%%%%%%%%%%%%%%%%%%%
\subsection{Modal Growth of Perturbations}
\frame
{
\hspace*{-0.75cm}
\Bmp{12cm}
\raggedright

\begin{itemize}
\item<1-> Numerical solution of the eigenvalue problem \blue{$\L \phi
    (\x) = \lambda \phi (\x)$}

 \medskip
\item<2->  Discretize $\L$ using the basis
  \footnotesize
\begin{equation*}
\begin{gathered}
\varphi_{j_1, j_2} = \frac{1}{\sqrt{2} \pi}\cos(j_1x + j_2y), \qquad 
\psi_{j_1, j_2} = -\frac{1}{\sqrt{2} \pi}\sin(j_1x + j_2y), \\
\qquad (j_1, j_2) \in \{(0, j_2), 1\leq j_2 \leq N\} \bigcup 
\{(j_1, j_2), 1\leq j_1 \leq N, -N\leq j_2\leq N\}
\end{gathered}
\end{equation*}
\normalsize
to obtain a sparse algebraic eigenvalue problem  \quad \blue{$\mathbf{L} \,
  \phi = \lambda \phi$}

\medskip
\item<3-> The spectrum of $\mathbf{L}$
\mbox{
  \visible<3->{\includegraphics[width=0.4\textwidth]{./Figs9/spectrum_TG}} \qquad
  \visible<4->{\includegraphics[width=0.4\textwidth]{./Figs9/eigs_Krylov}}}
\end{itemize}

     
\Emp
}


%%%%%%%%%%%%%%%%%%%%%%%%%%%%%%%%%%%%%%%%%%%%%%%%%%%%%%%%%%%%%%%%%%%%%%%%%%
%%%%%%%%%%%%%%%%%%%%%%%%%%%%%%%%%%%%%%%%%%%%%%%%%%%%%%%%%%%%%%%%%%%%%%%%%%
%%%%%%%%%%%%%%%%%%%%%%%%%%%%%%%%%%%%%%%%%%%%%%%%%%%%%%%%%%%%%%%%%%%%%%%%%%

\frame
{
\hspace*{-0.75cm}
\Bmp{12cm}
\raggedright

\begin{itemize}
\item<1-> The unstable eigenvalue: \quad \red{$\lambda_+^{3000} =
    0.1424 + 0.5875i$}

  \medskip
\item<2-> The corresponding eigenvectors $\Re\left[\phi_+^N\right]$ are
  $L^2$ integrable
\visible<2->{\mbox{
  \includegraphics[width=0.425\textwidth]{./Figs9/eigfun3000} \quad
  \includegraphics[width=0.425\textwidth]{./Figs9/vort_section}}
}

% \medskip
\vspace*{-0.25cm}
\item<3-> From the theorem of Lin (2004):
  \small \vspace*{-0.25cm}
  \begin{equation*}
\phi_+ \in W^{1,p^*}(\TT^2), \qquad \text{where} \qquad \blue{p^* = \frac{1}{1 -
    0.14} = 1.16}
\end{equation*}

\normalsize
\item<3-> From the Sobolev embedding  $W^{1,p^*} \hookrightarrow H^s$ with
  $1/p^* - 1/2 = 1/2 - s / 2$, we have
    \small \vspace*{-0.35cm}
  \begin{equation*}
\qquad \qquad \qquad \qquad \blue{\phi_+ \in H^{0.28}(\TT^2)}
\end{equation*}
  
\end{itemize}

     
\Emp
}

%%%%%%%%%%%%%%%%%%%%%%%%%%%%%%%%%%%%%%%%%%%%%%%%%%%%%%%%%%%%%%%%%%%%%%%%%%
%%%%%%%%%%%%%%%%%%%%%%%%%%%%%%%%%%%%%%%%%%%%%%%%%%%%%%%%%%%%%%%%%%%%%%%%%%
%%%%%%%%%%%%%%%%%%%%%%%%%%%%%%%%%%%%%%%%%%%%%%%%%%%%%%%%%%%%%%%%%%%%%%%%%%

\frame
{
\hspace*{-0.75cm}
\Bmp{12cm}
\raggedright

The Fourier spectrum of the eigenfunction $\phi$
\\ \vspace*{-0.4cm}
\begin{equation*}
e(k) := \frac{1}{2}\sum_{k \leq |\bds j| < k+1} |\widehat\phi_{\bds j}|^2, \qquad k \in \mbb N
\end{equation*}
\vspace*{-0.3cm}

\centering
\includegraphics[width=0.65\textwidth]{./Figs9/ori_eigfun_spec_s_0} \\
\small
$e(k) = \mathcal{O}\left(k^{-1}\right)$ implies $L^2$ summability;
$e(k) = \mathcal{O}\left(k^{-1.65}\right)$ means \blue{$\phi_+ \in H^{0.28}(\TT^2)$}

     
\Emp
}

%%%%%%%%%%%%%%%%%%%%%%%%%%%%%%%%%%%%%%%%%%%%%%%%%%%%%%%%%%%%%%%%%%%%%%%%%%
%%%%%%%%%%%%%%%%%%%%%%%%%%%%%%%%%%%%%%%%%%%%%%%%%%%%%%%%%%%%%%%%%%%%%%%%%%
%%%%%%%%%%%%%%%%%%%%%%%%%%%%%%%%%%%%%%%%%%%%%%%%%%%%%%%%%%%%%%%%%%%%%%%%%%

\frame
{
\hspace*{-0.75cm}
\Bmp{12cm}
\raggedright

Exponential growth of the perturbation
\bigskip

\centering
\mbox{
  \includegraphics[width=0.475\textwidth]{./Figs9/eig_linear_growth} \quad
  \includegraphics[width=0.475\textwidth]{./Figs9/eig_linear_growthrate}}
\\
\medskip
The actual growth rate is within 0.01\% equal $\Re(\lambda_+^{3000})$

\Emp
}


%%%%%%%%%%%%%%%%%%%%%%%%%%%%%%%%%%%%%%%%%%%%%%%%%%%%%%%%%%%%%%%%%%%%%%%%%%
%%%%%%%%%%%%%%%%%%%%%%%%%%%%%%%%%%%%%%%%%%%%%%%%%%%%%%%%%%%%%%%%%%%%%%%%%%
%%%%%%%%%%%%%%%%%%%%%%%%%%%%%%%%%%%%%%%%%%%%%%%%%%%%%%%%%%%%%%%%%%%%%%%%%%
\frame
{
\hspace*{-0.75cm}
\Bmp{12cm}
\raggedright
\centering

\mbox{
\includegraphics[width=0.45\textwidth]{./Figs9/growth_nonlinear}\qquad
\includegraphics[width=0.45\textwidth]{./Figs9/growth_rate_nonlinear}
}
\small
\medskip

The time-dependence of (left) the norm
{$\| \omega^{\delta}(t) - \omega_s\|_{H^{-1}}$} and (right) the growth
rate $(d/dt) \ln (\| \omega^{\delta}(t) - \omega_s\|_{H^{-1}})$ in the
solution of the nonlinear problem with the initial condition given in
terms of the eigenfunction $\phi_+^{200}$ as
$\omega_0 = \omega_s + \delta \cdot
\Re\left(\phi_+^{200}\right)/||\Re\left(\phi_+^{200}\right)||_{H^{-1}}$
with different indicated magnitudes $\delta$.

\Emp
}


%%%%%%%%%%%%%%%%%%%%%%%%%%%%%%%%%%%%%%%%%%%%%%%%%%%%%%%%%%%%%%%%%%%%%%%%%%
%%%%%%%%%%%%%%%%%%%%%%%%%%%%%%%%%%%%%%%%%%%%%%%%%%%%%%%%%%%%%%%%%%%%%%%%%%
%%%%%%%%%%%%%%%%%%%%%%%%%%%%%%%%%%%%%%%%%%%%%%%%%%%%%%%%%%%%%%%%%%%%%%%%%%
\subsection{Nonmodal Growth of Perturbations}
\frame
{
\hspace*{-0.75cm}
\Bmp{12cm}
\raggedright

\begin{itemize}
\item<1-> The growth rate of the eigenfunction with $\Re(\lambda_+)
  \approx 0.1424$ does not saturate the bound on the growth of the
  semigroup
\begin{equation*}
\blue{ \left \| e^{\L t} w_0 \right \|_{H^{m}} \le \left \| w_0 \right \|_{H^{m}} e^{\mumax t}, \qquad \mumax = 1, \quad m = \pm 1}
\end{equation*}

\bigskip
\item<2-> How to find initial data $w_0$ realizing this growth?

 \bigskip
\item<3->
  Given $T > 0$, define the objective functional $J_m \; : \;
  {H_0^m(\TT^2)} \rightarrow \RR$
\begin{equation*}
\blue{  J_m (w_0) = \left|\left|e^{T\L} w_0 \right|\right|^2_{H^m},
\qquad m = \pm 1}
\end{equation*}
  

\end{itemize}
\Emp
}

%%%%%%%%%%%%%%%%%%%%%%%%%%%%%%%%%%%%%%%%%%%%%%%%%%%%%%%%%%%%%%%%%%%%%%%%%%
%%%%%%%%%%%%%%%%%%%%%%%%%%%%%%%%%%%%%%%%%%%%%%%%%%%%%%%%%%%%%%%%%%%%%%%%%%
%%%%%%%%%%%%%%%%%%%%%%%%%%%%%%%%%%%%%%%%%%%%%%%%%%%%%%%%%%%%%%%%%%%%%%%%%%
\frame
{
\hspace*{-0.75cm}
\Bmp{12cm}
\raggedright

\begin{itemize}
\item<1-> []
\begin{problem}[4]
\vspace*{-0.5cm}
\blue{\small
\begin{align*}
& \max_{w_0 \in \mc M} J_m (w_0), \quad \text{where}  \ m = \pm 1  \ \text{and} \\
& \qquad \mc M := \{w_0 \in H^m_0(\TT^2):  \|w_0\|_{H^m} =   1\}. \\ \\
& \text{subject to:} \quad  \left\{
\begin{aligned}
&\partial_tw - \L w = 0,  && \text{in} \ \TT^2\times(0,T] \\ 
&w = w_0                 &&  \text{in} \ \TT^2 \ \text{at} \ t=0 \\
&\textrm{Periodic Boundary Conditions}                  &&  
\end{aligned}
%\label{eq:NS}
\right.
\end{align*}
}
\end{problem}

\item<2->
  Problem solved with a discretized gradient flow, as before
\end{itemize}
\Emp
}


%%%%%%%%%%%%%%%%%%%%%%%%%%%%%%%%%%%%%%%%%%%%%%%%%%%%%%%%%%%%%%%%%%%%%%%%%%
%%%%%%%%%%%%%%%%%%%%%%%%%%%%%%%%%%%%%%%%%%%%%%%%%%%%%%%%%%%%%%%%%%%%%%%%%%
%%%%%%%%%%%%%%%%%%%%%%%%%%%%%%%%%%%%%%%%%%%%%%%%%%%%%%%%%%%%%%%%%%%%%%%%%%
\frame
{
\hspace*{-0.75cm}
\Bmp{12cm}
\raggedright
\centering

\begin{figure}
\mbox{\subfigure[$N=128$]{\includegraphics[width=0.325\textwidth]{./Figs9/s_1_128}}\qquad
\subfigure[$N=256$]{\includegraphics[width=0.325\textwidth]{./Figs9/s_1_256}}}
\mbox{\subfigure[$N=512$]{\includegraphics[width=0.325\textwidth]{./Figs9/s_1_512}}\qquad
\subfigure[$N=1024$]{\includegraphics[width=0.325\textwidth]{./Figs9/s_1_1024}}}
\end{figure}

\small\vspace*{-0.3cm}
\blue{[$m=1$]}  The optimal initial conditions $\widetilde{w}_0^N$ obtained using
different resolutions $N$.

\Emp
}


%%%%%%%%%%%%%%%%%%%%%%%%%%%%%%%%%%%%%%%%%%%%%%%%%%%%%%%%%%%%%%%%%%%%%%%%%%
%%%%%%%%%%%%%%%%%%%%%%%%%%%%%%%%%%%%%%%%%%%%%%%%%%%%%%%%%%%%%%%%%%%%%%%%%%
%%%%%%%%%%%%%%%%%%%%%%%%%%%%%%%%%%%%%%%%%%%%%%%%%%%%%%%%%%%%%%%%%%%%%%%%%%
\frame
{
\hspace*{-0.75cm}
\Bmp{12cm}
\raggedright

\mbox{
\includegraphics[width=0.475\textwidth]{./Figs9/s_1_growth}
\includegraphics[width=0.475\textwidth]{./Figs9/approx_s_1_crlt}}
\centering
\small
\blue{[$m=1$]} (left) The growth rate
  $(d/dt) \ln (\|w^N(t)\|_{H^{1}})$ versus $t$ for
  the optimal initial conditions $\widetilde{w}_0^N$ obtained using increasing
  spatial resolutions {$N^2$}. \\
  (right) The time-dependence of the
  autocorrelation function
  \begin{equation*}
\C^1_{\tau}(t) := \frac{\left\langle w(\tau), w(t) \right\rangle_{H^1}}{\| w(\tau) \|_{H^1} \| w(t) \|_{H^1}}, 
\qquad t, \tau \ge 0
\end{equation*}
  corresponding to
  $N = 1024$ and $\tau = 0, 0.25, 0.5, 0.75, 1$. \\
  Similar results obtained for $m = -1$.
\Emp
}



%%%%%%%%%%%%%%%%%%%%%%%%%%%%%%%%%%%%%%%%%%%%%%%%%%%%%%%%%%%%%%%%%%%%%%%%%%
%%%%%%%%%%%%%%%%%%%%%%%%%%%%%%%%%%%%%%%%%%%%%%%%%%%%%%%%%%%%%%%%%%%%%%%%%%
%%%%%%%%%%%%%%%%%%%%%%%%%%%%%%%%%%%%%%%%%%%%%%%%%%%%%%%%%%%%%%%%%%%%%%%%%%
\frame
{
\hspace*{-0.75cm}
\Bmp{12cm}
\raggedright
\centering

\movie[width=0.9\textwidth,loop,showcontrols]{\includegraphics[width=0.9\textwidth]{./Figs9/s_1_1024}}{../Euler2D_stability/Movies/Movie1.mov}

%\movie[width=0.8\textwidth,loop,showcontrols]{\includegraphics[width=0.8\textwidth]{./Figs9/s_1_1024}}{../Euler2D_stability/Movies/movie1.mp4 }

\Emp
}



%%%%%%%%%%%%%%%%%%%%%%%%%%%%%%%%%%%%%%%%%%%%%%%%%%%%%%%%%%%%%%%%%%%%%%%%%%
%%%%%%%%%%%%%%%%%%%%%%%%%%%%%%%%%%%%%%%%%%%%%%%%%%%%%%%%%%%%%%%%%%%%%%%%%%
%%%%%%%%%%%%%%%%%%%%%%%%%%%%%%%%%%%%%%%%%%%%%%%%%%%%%%%%%%%%%%%%%%%%%%%%%%
\frame
{
\frametitle{Conclusions}
\vspace{-0.1cm}
%\vspace*{0.0cm}
\hspace*{-1.0cm}
\Bmp{12cm}
\raggedright
\small
\begin{itemize}
\item<1-> The Taylor-Green vortex is linearly unstable
  \begin{itemize}
\item<2-> it has a eigenmode given by a distribution with the Sobolev
  regularity consistent with Lin (2004)
  \smallskip 
\item<3-> it also exhibits a nonmodal growth mechanism involving a
  continuous family of uncorrelated functions which saturates the
  spectral bound
 \begin{itemize}
 \item<4-> intrinsically related to the $\infty$-dimensional nature of the operator $\L$ 
 \smallskip
 \item<5-> this mechanism is consistent with the WKB analysis based on
   the bicharacteristic equation
\end{itemize}
\end{itemize}

\bigskip
\item<6-> The unstable eigenmode leads to a nonlinear instability
%\begin{itemize}
%\item<7->
%using the Sobolev embedding $W^{1,p^*} \hookrightarrow L^q$, 
%$1/p^* - 1/2 = 1/q$, \\ we have \blue{$\phi_+ \in L^q(\TT)$} with \blue{$q = 2.76$}
%\smallskip
%\item<8->
%thus, the initial condition $\omega_s + \delta \Re\left(\phi_+\right)$ is {\em not} 
%in the  Yudovich class $L^1(\TT^2) \bigcap L^\infty(\TT^2)$ \qquad 
%\visible<8->{$\Longrightarrow$ \quad \blue{possible nonuniqueness?}}
%\end{itemize}
  
\bigskip  
\item<7-> \red{\sc Open Question:} stability of equilibria in 3D
\end{itemize}
\Emp
}



\end{document}

%%%%%%%%%%%%%%%%%%%%%%%%%%%%%%%%%%%%%%%%%%%%%%%%%%%%%%%%%%%%%%%%%%%%%%%%%%
%%%%%%%%%%%%%%%%%%%%%%%%%%%%%%%%%%%%%%%%%%%%%%%%%%%%%%%%%%%%%%%%%%%%%%%%%%
%%%%%%%%%%%%%%%%%%%%%%%%%%%%%%%%%%%%%%%%%%%%%%%%%%%%%%%%%%%%%%%%%%%%%%%%%%

\frame
{
\frametitle{References}
%\vspace{-0.15cm}
\vspace*{0.0cm}
\hspace*{-1.0cm}
\Bmp{12cm}
\raggedright
\scriptsize
\begin{itemize}
\item
L.~Lu and C.~R.~Doering, ``Limits on Enstrophy Growth for Solutions of
the Three-dimensional Navier-Stokes Equations'' {\em Indiana
University Mathematics Journal} {\bf{57}}, 2693--2727, 2008.

\vspace*{-0.3cm}
\begin{center}
\rule{0.8\textwidth}{0.025cm}
\end{center}
\vspace*{-0.3cm}

\smallskip
\item
D.~Ayala and B.~Protas, ``On Maximum Enstrophy Growth in a
Hydrodynamic System'', {\em Physica D} {\bf{240}}, 1553--1563, 2011.

\smallskip
\item D.~Ayala and B.~Protas, ``Maximum Palinstrophy Growth in 2D
  Incompressible Flows: Instantaneous Case'', {\em Journal of Fluid
    Mechanics} {\bf{742}} 340--367, 2014.

\smallskip
\item D.~Ayala and B.~Protas, ``Vortices, Maximum Growth and the
  Problem of Finite-Time Singularity Formation'', {\em Fluid Dynamics
    Research}, {\bf{46}}, 031404, 2014.

\smallskip
\item D.~Ayala and B.~Protas, ``Extreme Vortex States and the Growth
  of Enstrophy in 3D Incompressible Flows'', {\em Journal of Fluid
    Mechanics} {\bf{818}}, 772--806, 2017.

\smallskip
\item D.~Po\c{c}as and B. Protas, ``Transient Growth in Stochastic
  Burgers Flows'', {\em Discrete and Continuous Dynamical Systems ---
    B} {\bf{23}}, 2371--2391, 2018.

\smallskip
\item D.~Yun and B.~Protas, ``Maximum Rate of Growth of Enstrophy in
  Solutions of the Fractional Burgers Equation'', {\em Journal of
    Nonlinear Science} {\bf{28}}, 395-422, 2018.

\smallskip
\item D.~Kang, D.~Yun and B.~Protas, ``Maximum Amplification of
  Enstrophy in 3D Navier-Stokes Flows'', (see {\tt arXiv:1909.00041}),
  2019.
\end{itemize}
\Emp
}



\end{document}


% TODO 
% first slide in ``Results'' - add arrows in the figures to indicate 
% trends w.r.t T and E0
% first slide in ``Envelopes'' - add line indicating the power law fit
